\documentclass[10pt, letterpaper]{article}

% Inhaltsverzeichnis für Pakettypen (nur für Übersicht im Header, wird nicht im Dokument angezeigt)
% 1. Seitenlayout und Ränder
% 2. Sprache und Zeichensatz
% 3. Mathematik und Theorem-Umgebungen
% 4. Eigene Makros
% 5. Diagramme und Grafiken
% 6. Tabellen und Aufzählungen
% 7. Inhaltsverzeichnis
% 8. Abschnittsüberschriften
% 9. Abstrakt-Umgebung
% 10. Todos/Notizen
% 11. Rahmen/Box-Umgebungen
% 12. Python-Integration
% 13. Literaturverwaltung
% 14. Hyperlinks
% 15. Absatzeinstellungen
% 16. Umgebungen
% 17  Graphik
% 18  Extra
% 00. Titel und Autor

\usepackage{slashed}

% --- 1. Seitenlayout und Ränder ---
\usepackage[margin=3cm]{geometry}

% --- 2. Sprache und Zeichensatz ---
\usepackage[english]{babel}
\usepackage[T1]{fontenc}
\usepackage[utf8]{inputenc}

% --- 3. Mathematik und Theorem-Umgebungen ---
\usepackage{amsmath, amssymb, amsthm}
\usepackage{mathrsfs}
\DeclareMathOperator{\WF}{WF}

% --- 4. Eigene Makros ---
\usepackage{xcolor}
\newcommand{\SKP}{\langle\cdot,\cdot\rangle}
\newcommand{\id}{\operatorname{id}}
\newcommand{\sgn}{\operatorname{sgn}}
\newcommand{\Ad}{\operatorname{Ad}}
\newcommand{\ad}{\operatorname{ad}}
\newcommand{\K}{\mathbb{K}}
\newcommand{\R}{\mathbb{R}}
\newcommand{\N}{\mathbb{N}}
\newcommand{\Q}{\mathbb{Q}}
\newcommand{\Z}{\mathbb{Z}}
\newcommand{\C}{\mathbb{C}}
\newcommand{\QU}{\mathbb{H}}
\newcommand{\Cinf}{C^{\infty}}
\newcommand{\Mod}{\mathrm{Mod}}
\newcommand{\Alg}{\mathrm{Alg}}
\newcommand{\Diff}{\mathrm{Diff}}
\newcommand{\Iso}{\mathrm{Iso}}
\newcommand{\Aut}{\mathrm{Aut}}
\newcommand{\End}{\mathrm{End}}
\newcommand{\Hom}{\mathrm{Hom}}
\newcommand{\Cl}{\mathrm{Cl}}
\newcommand{\entwurf}[1]{\textcolor{red}{#1}}
\newcommand{\GL}{\mathrm{GL}}
\newcommand{\Mat}{\mathrm{Mat}}
\newcommand{\SL}{\mathrm{SL}}
\newcommand{\SO}{\mathrm{SO}}
\newcommand{\SU}{\mathrm{SU}}
\newcommand{\Sp}{\mathrm{Sp}}
\newcommand{\Spin}{\mathrm{Spin}}
\newcommand{\Pin}{\mathrm{Pin}}
\newcommand{\Lor}{\mathrm{Lor}}
\newcommand{\Ogroup}{\mathrm{O}}
\newcommand{\U}{\mathrm{U}}
\newcommand{\CP}{\mathbb{C} P}
\newcommand{\RP}{\mathbb{R} P}

% --- 5. Diagramme und Grafiken ---
\usepackage{graphicx}
\graphicspath{{../../Library/Images/}}
\usepackage[export]{adjustbox}
\usepackage{tikz}
\usetikzlibrary{decorations.pathreplacing, arrows.meta, positioning}
\usepackage{tikz-cd}

% --- 6. Tabellen und Aufzählungen ---
\usepackage{enumitem}
\setlist[itemize]{left=0.5cm}

\newenvironment{romanenum}[1][]
  {%
    \ifx&#1&
    \else
      \textbf{#1}\quad
    \fi
    \begin{enumerate}[label=\roman*)]
  }
  {%
    \end{enumerate}%
  }

% --- 7. Inhaltsverzeichnis ---
\usepackage{tocloft}
\renewcommand{\cftsecfont}{\footnotesize}
\renewcommand{\cftsubsecfont}{\footnotesize}
\renewcommand{\cftsubsubsecfont}{\footnotesize}
\renewcommand{\cftsecpagefont}{\footnotesize}
\renewcommand{\cftsubsecpagefont}{\footnotesize}
\renewcommand{\cftsubsubsecpagefont}{\footnotesize}
\usepackage{etoc}

% --- 8. Abschnittsüberschriften ---
\usepackage{titlesec}
\titleformat{\section}{\normalfont\large\bfseries}{\thesection}{1em}{}
\titleformat{\subsection}{\normalfont\normalsize\bfseries}{\thesubsection}{0.5em}{}
\titleformat{\subsubsection}{\normalfont\normalsize\bfseries}{\thesubsubsection}{0.5em}{}
\setcounter{secnumdepth}{4}

% --- 9. Abstrakt-Umgebung ---
\usepackage{changepage}
\renewenvironment{abstract}
  {
    \begin{adjustwidth}{1.5cm}{1.5cm}
    \small
    \textsc{Abstract. –}%
  }
  {
    \end{adjustwidth}
  }

% --- 10. Todos/Notizen ---
\usepackage{todonotes}
% Hellblau definieren
\definecolor{hellblau}{rgb}{0.3,0.6,1}
\newenvironment{note}{\par\begingroup\color{hellblau}\itshape}{\par\endgroup}

% --- 11. Rahmen/Box-Umgebungen ---
\usepackage{mdframed}
\usepackage{tcolorbox}
\colorlet{shadecolor}{gray!25}

\newenvironment{customTheorem}
  {\vspace{10pt}%
   \begin{mdframed}[
     backgroundcolor=gray!20,
     linewidth=0pt,
     innertopmargin=10pt,
     innerbottommargin=10pt,
     skipabove=\dimexpr\topsep+\ht\strutbox\relax,
     skipbelow=\topsep,
   ]}
  {\end{mdframed}
   \vspace{10pt}%
  }

% --- 12. Python-Integration ---
% (Deaktiviert in dieser Version, aktiviere bei Bedarf)
% \usepackage{pythontex}
% \usepackage[makestderr]{pythontex}

% --- 13. Literaturverwaltung ---
\usepackage{csquotes}
\usepackage[backend=biber, style=alphabetic, citestyle=alphabetic]{biblatex}
\addbibresource{bibliography.bib}

% --- 14. Hyperlinks ---
\usepackage{hyperref}
\hypersetup{
  colorlinks   = true,
  urlcolor     = blue,
  linkcolor    = blue,
  citecolor    = blue,
  frenchlinks  = true
}

% --- 15. Absatzeinstellungen ---
\usepackage[parfill]{parskip}
\sloppy

% --- 16. Umgebungen ---
\usepackage{thmtools}

\newcommand{\CustomHeading}[3]{%
  \par\medskip\noindent%
  \textbf{#1 #2} \textnormal{(#3)}.\enskip%
}

\newenvironment{DEF}[2]{\CustomHeading{Definition}{#1}{#2}}{}
\newenvironment{PROP}[2]{\CustomHeading{Proposition}{#1}{#2}}{}
\newenvironment{THEO}[2]{\CustomHeading{Theorem}{#1}{#2}}{}
\newenvironment{LEM}[2]{\CustomHeading{Lemma}{#1}{#2}}{}
\newenvironment{KORO}[2]{\CustomHeading{Corollar}{#1}{#2}}{}
\newenvironment{REM}[2]{\CustomHeading{Remark}{#1}{#2}}{}
\newenvironment{EXA}[2]{\CustomHeading{Example}{#1}{#2}}{}
\newenvironment{STUD}[2]{\CustomHeading{Study}{#1}{#2}}{}
\newenvironment{CONC}[2]{\CustomHeading{Concept}{#1}{#2}}{}
\newenvironment{OTH}[2]{\CustomHeading{Other}{#1}{#2}}{}
\newenvironment{EXE}[2]{\CustomHeading{Exercise}{#1}{#2}}{}
\newenvironment{MOT}[2]{\CustomHeading{Motivation}{#1}{#2}}{}
\newenvironment{PROOF}[2]{\CustomHeading{Proof}{#1}{#2}}{}

% --- Unit Umgebung für Source-Inhalte ---
\usepackage{mdframed}
\newmdenv[
  linewidth=1pt,
  topline=false,
  bottomline=false,
  rightline=false,
  leftmargin=0cm,
  rightmargin=0cm,
  skipabove=10pt,
  skipbelow=10pt,
  innertopmargin=0.5\baselineskip,
  innerbottommargin=0.5\baselineskip,
  backgroundcolor=gray!10,
  linecolor=gray
]{unitbox}

\newenvironment{unit}[1]
  {\begin{unitbox}\textbf{Unit #1}\par\smallskip}
  {\end{unitbox}}

% --- 17. ... ---

% --- 18. Extras ---
\usepackage{stmaryrd}
\usepackage{bbold}  % falls du \mathbb{1} nutzen willst


% --- 19. Inhaltsverzeichnis ---
\usepackage{tocloft}

% Abstand zwischen Einträgen
\setlength{\cftbeforesecskip}{2pt}      % Sections
\setlength{\cftbeforesubsecskip}{1pt}   % Subsections
\setlength{\cftbeforesubsubsecskip}{0pt}
\renewcommand{\cfttoctitlefont}{\large\bfseries}



% --- 00. Titel und Autor ---
\title{Einführung in Riemannsche Spin-Geometrie}
\author{Axel Jaschik}
\date{\today}




\begin{document}



\begin{center}
    \Large{Einführung in Riemannsche Spin-Geometrie}\\[10pt]
\end{center}
\vspace{0.05cm}

\begin{center}
     {Herrn Prof. Dr. Bernig: \textit{Riemannsche Geometrie und Lie Gruppen} (Blockseminar WS25/26)}\\
\end{center}


\begin{center}
    \footnotesize{Axel Tim Jaschik \qquad s0852124@stud.uni-frankfurt.de\qquad Matrik.Nr.: 6641832}
\end{center}

\rule{\textwidth}{0.5pt}
\begin{abstract}
Diese Ausarbeitung gibt eine erste, möglichst geschlossene Einführung in die Riemannsche Spin-Geometrie mit Schwerpunkt auf der differentialgeometrischen Motivierung und Konstruktion des Dirac-Operators auf Spinorbündeln. Ausgehend von einer natürlichen Fragestellung über Differentialoperatoren vom Laplace-Typ auf Vektorbündeln wird ein Ansatz entwickelt, in welchem sich zeigt, warum \emph{Clifford-Strukturen}, \emph{Spin-Strukturen} und \emph{Spinorbündel} als die natürliche algebraische bzw. geometrischen Strukturen aufkommen. Es wird die \emph{Clifford-Multiplikation} $c$ und der \emph{Spinorzusammenhang} $\nabla^\Sigma$ auf dem Spinorbündel $\Sigma M$ konstruiert, aus welchen dann der \emph{Dirac-Operator} zusammengesetzt wird
$$D=c\circ \nabla^\Sigma.$$ 
Abschließend wird die \emph{Lichnerowicz-Formel} für den \emph{Dirac-Laplace-Operator} $D^2$ sowie ein \emph{Vanishing-Theorem} für harmonische Spinorfelder für kompakte Riemannsche Spin-Mannigfaltigkeiten ohne Rand mit positiver Skalarkrümmung bewiesen.
\end{abstract}
\rule{\textwidth}{0.5pt}



\tableofcontents


\vspace{0.8cm}


\section{Problemstellung}

Diese Ausarbeitung soll eine Einführung in die sogenannte \emph{Riemannsche Spin-Geometrie} geben. Da allein die Etablierung der Grundobjekte der Spin-Geometrie einiges an konzeptueller Arbeit erfordert, wird es notwendig sein, an einigen Stellen Details kurz zu referieren und auf die Literatur zu verweisen, um abschließend noch erste interessante Resultate erreichen zu können. 


Wir betrachten eine orientierte Riemannsche Mannigfaltigkeit
\[
(M^n,g,\mathrm{Orie}),
\qquad n=2m,\quad m\in\mathbb N,
\]
gerader Dimension. 
(Auf die Rolle dieser Dimensionsannahme gehen wir später noch genauer ein.) Der Zugang zur Spin-Geometrie soll über folgende natürliche geometrische Problemstellung erfolgen.

\begin{MOT}{}{Natürlicher Differentialoperator zwischen Vektorbündeln}
  Gibt es ein Tripel 
  $$\left( E,\nabla^E, D \right)$$
  von $E=(E,\pi,M,V)$ einem Vektorbündel mit typischer Faser $V$ einem endlichdimensionalen $\mathbb{K}$-Vektorraum und $\nabla^E$ eine Zusammenhang auf $E$ sowie einem Differentialoperator $D\in \mathscr{D}_1(\Gamma(E),\Gamma(E)))$, sodass folgende Eigenschaften durch das Tripel erfüllt werden:
  \begin{enumerate}
    \item \emph{Natürlichkeit:}
    Das Bündel $E$, der Zusammenhang $\nabla^E$ und der Operator $D$ sollen
    in kanonischer Weise aus den Riemannschen Daten $(M,g,\mathrm{Orie})$
    hervorgehen.
    \item \emph{Laplace-Typ-Bedingung für $D^2$:}
    Bezüglich des Zusammenhangs $\nabla^E$ auf $E$ soll in jedem lokalen orthonormalen Rahmen $(e_1,\dots,e_n)$ folgende Form gelten:
    \[
    D^2
    =
    -\sum_{i=1}^n\Bigl(
    \nabla^E_{e_i}\nabla^E_{e_i}
    -\nabla^E_{\nabla^{LC}_{e_i}e_i}
    \Bigr)
    +0.o.t.
    \]
  \end{enumerate}
\end{MOT}
\bigskip 

Die Riemannsche Spin-Geometrie liefert eine Antwort auf diese Frage und entwickelt eine umfassende Theorie auf geeigneten Tripeln von Vektorbündeln, Zusammenhängen und Differentialoperatoren mit weitreichenden Resultaten an der Schnittstelle von Differentialgeometrie und Topologie, aber auch in der Mathematischen Physik (Pseudoriemannsche Spin-Geometrie).


Diese Ausarbeitung widmet sich der Entwicklung des kanonischen Tripels, mit welchem die Riemannsche Spin-Geometrie klassischerweise ansetzt: Das \emph{Spinorbündel}, der \emph{Spinorzusammenhang} und der \emph{Dirac-Operator}. Dabei orientieren wir uns primär an Bärs Lecture Notes \emph{Spin Geometry} \parencite[]{BaerSpin}.


\section{Natürliche Clifford-Strukturen}


\subsection{Laplace-Typ-Bedingung erzwingt Realisierung von Clifford-Relation durch Koeffizientenabbildungen in lokaler Darstellung}

Lassen wir zunächst die Frage nach dem geeigneten Vektorbündel außen vor und betrachten genauer die lokale Version der Laplace-Typ-Bedingung, so finden wir direkt die charakteristische Eigenschaft, die Differentialoperatoren der Spin-Geometrie auszeichnet.



Sei \(e_1,\dots,e_n\) ein lokaler orthonormaler Rahmen von \(TM\). Angenommen wir haben einen natürlichen Differentialoperator $D$ und Zusammenhang $\nabla^E$. Da \(D\) ein Differentialoperator erster Ordnung ist, kann er lokal in der Form
\[
D=\sum_{i=1}^n a_i\nabla^E_{e_i}+ \text{0.o.t.}
\]
geschrieben werden, wobei wir $a_i$ zunächst als nicht weiter spezifizierte Koeffizientenabbildungen auffassen. Für die Analyse der Laplace-Typ-Bedingung von \(D^2\) ist zunächst nur der Hauptterm relevant
\[
D\sim \sum_{i=1}^n a_i\nabla^E_{e_i}.
\]
Beim Quadrieren ergibt sich im Term zweiter Ordnung
\[
D^2\sim \sum_{i,j=1}^n a_i a_j\,\nabla^E_{e_i}\nabla^E_{e_j}.
\]
Symmetrisieren liefert
\[
\sum_{i=1}^n a_i^2\,\nabla^E_{e_i}\nabla^E_{e_i}
+\frac12\sum_{i\neq j}(a_i a_j+a_j a_i)\,\nabla^E_{e_i}\nabla^E_{e_j}.
\]
Soll \(D^2\) vom Laplace-Typ sein, so muss der zweite Ordnungsterm die geforderten Form haben. Aus Koeffizientenvergleich folgen notwendig die Relationen
\[
a_i^2=-1,
\qquad
a_i a_j+a_j a_i=0 \quad i\neq j,
\]
also zusammengefasst
\[
a_i a_j+a_j a_i=-2\delta_{ij}.
\]
Das hat gravierende Konsequenzen, denn damit zeigt sich die notwendig nicht-triviale Natur der Koeffizientenabbildungen von natürlichen Differentialoperatoren. Da es sich insbesondere bei den Koeffizientenabbildungen nicht mehr um skalare Funktionen handeln kann, die punktweise einfach skalar auf $\nabla^E_{X}$ wirken, müssen wir hier eine komplexere Form von Wirkung vermuten:
$$a_i(x)\in\End(E_x) \quad \text{ wirken auf }\quad  \nabla^E_Xs(x)\in E_x,$$
wobei wir nun die $a_i$ Abbildungen punktweise als \emph{Koeffizientenendomorphismen} auffassen, welche auf $\nabla^E_X s$ angewandt werden und unter sich eine Clifford-Relation realisieren müssen. Die Koeffizientenendomorphismen verhalten sich also wie eine \emph{Realisierung der formalen Clifford-Erzeuger auf dem Vektorraum $E_x$}.




\subsection{Allgemeine Clifford-Strukturen}

Um über Clifford-Strukturen in einem allgemeinen Kontext sprechen zu können, müssen wir einige Begriffe einführen.

\begin{DEF}{}{Clifford-Algebra von \((V,\beta)\)}
Sei \(\mathbb{K}=\mathbb{R}\) oder \(\mathbb{C}\), und sei \(V\) ein endlichdimensionaler \(\mathbb{K}\)-Vektorraum, ausgestattet mit einer symmetrischen Bilinearform
\[
\beta:V\times V\to \mathbb{K}.
\]

\begin{enumerate}
\item Eine \emph{Clifford-Algebra} zu \((V,\beta)\) ist eine unitale \(\mathbb{K}\)-Algebra
\(\Cl(V,\beta)\) zusammen mit einer linearen Abbildung
\[
\iota:V\to \Cl(V,\beta),
\]
so dass die folgenden Eigenschaften erfüllt sind:

\begin{enumerate}
    \item Für alle \(v\in V\) gilt
    \[
    \iota(v)^2=-\beta(v,v)\cdot 1.
    \]
    Äquivalent dazu gilt für alle \(v,w\in V\):
    \[
    \iota(v)\iota(w)+\iota(w)\iota(v)=-2\beta(v,w)\cdot 1.
    \]

    \item Das Paar \((\Cl(V,\beta),\iota)\) ist universal bezüglich der Clifford-Relation, d.h.:
    Ist \(A\) eine unitale \(\mathbb{K}\)-Algebra und
    \[
    \iota':V\to A
    \]
    eine lineare Abbildung mit
    \[
    \iota'(v)^2=-\beta(v,v)\cdot 1
    \qquad\text{für alle }v\in V,
    \]
    so existiert genau ein Algebra-Homomorphismus
    \[
    \varphi:\Cl(V,\beta)\to A,
    \]
    so dass das Diagramm
    \[
    \begin{tikzcd}
    V \arrow[r, "\iota"] \arrow[dr, "\iota'"'] & \Cl(V,\beta) \arrow[d, "\varphi"] \\
    & A
    \end{tikzcd}
    \]
    kommutiert.
\end{enumerate}

\item Falls der Vektorraum $V$ reell ist, ist die komplexifizierte Clifford-Algebra von $\Cl_n(V,\beta)$ definiert durch 
  \[
  \C l(V,g):=\Cl(V,g)\otimes_{\R}\C.
  \]
\end{enumerate}
\end{DEF}


\bigskip




Wir halten wichtige Eigenschaften der Clifford-Algebra fest:

\begin{REM}{}{Eigenschaften der Clifford-Algebra}
\begin{enumerate}
  \item \emph{Basis von $\Cl(V,\beta)$}: Sei $V$ ein $\K$-Vektorraum mit Basis $b_1, \ldots, b_n \in V$. Dann ist 
    $$
    \left\{b_{i_1} \cdot \ldots \cdot b_{i_k}\right\}_{\substack{k=0,1, \ldots n \\ 1 \leq i_1<i_2<\ldots<i_k \leq n}}
    $$
    eine Basis der Clifford-Algebra mit Dimension $\operatorname{dim}_{\mathbb{K}} \mathrm{Cl}(V, \beta)=2^{\operatorname{dim}_{\mathbb{K}} V}$.
  \item \emph{$\Z_2$-Graduierung der Clifford Algebra und kanonischer Automorphismus}: Sei $\imath: V \rightarrow \mathrm{Cl}(V, \beta)$ die Standardeinbettung und $\imath^{\prime}: V \rightarrow \mathrm{Cl}(V, \beta)$ definiert durch
    $$
    \imath^{\prime}(v)=-\imath(v), \quad  v \in V .
    $$
    Da $\imath^{\prime}(v)^2=\imath(v)^2=-\beta(v, v) \cdot 1$ gilt, existiert wegen universeller Eigenschaft der eindeutige \emph{kanonische Algebren-Automorphismus} $\alpha: \mathrm{Cl}(V, \beta) \rightarrow \mathrm{Cl}(V, \beta)$, so dass das Diagramm
    \[
    \begin{tikzcd}
    V \arrow[r, "\iota"] \arrow[dr, "\iota'"'] & \Cl(V,\beta) \arrow[d, "\alpha"] \\
    & \Cl(V,\beta)
    \end{tikzcd}
    \]
    kommutiert. Durch Identifikation $\imath(V) \subset \mathrm{Cl}(V, \beta)$ mit $V$ haben wir $\left.\alpha\right|_V=-\mathrm{id}$. Das induziert eine Dekomposition 
    $$
    \mathrm{Cl}(V, \beta)=\mathrm{Cl}^0(V, \beta) \oplus \mathrm{Cl}^1(V, \beta),
    $$
    mit
    $$
    \begin{aligned}
    & \mathrm{Cl}^0(V, \beta)=+1-\text {Eigenraum von } \alpha \\
    & \mathrm{Cl}^1(V, \beta)=-1-\text {Eigenraum von } \alpha.
    \end{aligned}
    $$
\end{enumerate}
\end{REM}




\begin{DEF}{}{Clifford-Modul}
Sei $\C l(V,\beta)$ eine komplexifizierte Clifford-Algebra einer reellen Clifford-Algebra. Ein \emph{(komplexes linkes) Clifford-Modul} über $\C l(V,\beta)$ ist ein komplexer Vektorraum $W$
zusammen mit einer Algebrenabbildung
\[
\rho:\C l(V,\beta)\longrightarrow \End_{\C}(W).
\]
Die dadurch induzierte Wirkung
\[
\C l(V,\beta)\times W\longrightarrow W,\qquad
(z,\psi)\longmapsto z\cdot \psi:=\rho(z)(\psi)
\]
heißt die \emph{Wirkung der Clifford-Algebra auf $W$}.

Die \emph{Clifford-Multiplikation} auf einem komplexen Clifford-Modul $W$ über
$\C l(V,\beta)$ ist die Einschränkung (bzgl. kanonischer Einbettung $V\subset \C l(V,\beta)$)
\[
c:=\rho|_V:V\to \End_{\C}(W)
\]
und man schreibt
\[
c(v)\psi=v\cdot\psi \qquad v\in V,\ \psi\in W.
\]
Die Clifford-Relation lautet dann
\[
c(v)c(w)+c(w)c(v)=-2\,g(v,w)\,\id_{W}.
\]
Wegen der Clifford-Relation bestimmt $c$ mittels der universellen Eigenschaft von $\C l(V, \beta)$ eindeutig die Darstellung $\rho$.
\end{DEF}


\subsection{Geometrische Clifford-Strukturen über $(T_xM,g_x)$}


\begin{DEF}{}{Clifford-Struktur über $(T_xM,g_x)$}
Für $x\in M$ bezeichne 
$$\Cl(T_xM,g_x)$$ 
die Clifford-Algebra des metrischen Vektorraums $(T_xM,g_x)$.

Wir vermuten also, dass die Clifford-Strukturen, welche in der lokalen Darstellung faserweise auf $E_x$ aufgetaucht sind, von einer allgemeineren geometrische Clifford-Struktur $\C l(T_xM,g_x)$ durch 
\[
\begin{tikzcd}
T_xM \arrow[r,"\iota_x"] \arrow[dr,"c_x"'] & \C l(T_xM,g_x) \arrow[d,"\widetilde c_x"] \\
& \End(E_x)
\end{tikzcd}
\]
herkommen, wobei $c_x$ die allgemeine Clifford-Relation realisiert und die Einschränkung von $\widetilde{c}_x$ auf kanonisch eingebettete $T_xM$ wieder $c_x$ entspricht.

Wir würden die konkreten Koeffizientenendomorphismen und deren Clifford-Relationen untereinander erhalten bzgl. des lokalen orthonormalen Rahmen \(e_1,\dots,e_n\) durch
$$c_x\left(\sum_{i=1}^n v^i e_i(x)\right):=\sum_{i=1}^n v^i a_i(x)$$
und können die Clifford-Relation einsehen:
$$\begin{aligned} c_x(v) c_x(w)+c_x(w) c_x(v) & =\sum_{i, j} v^i w^j\left(a_i a_j+a_j a_i\right) \\ & =\sum_{i, j} v^i w^j\left(-2 \delta_{i j}\right) \mathrm{id}_{E_x} \\ & =-2 \sum_i v^i w^i \mathrm{id}_{E_x} \\ & =-2 g_x(v, w) \mathrm{id}_{E_x},\end{aligned}$$
\end{DEF}



\subsection{Ansatz für natürliches Vektorbündel mit Clifford-Struktur}

Wir machen jetzt einen Ansatz für die Konstruktion eines natürlichen Vektorbündels, auf welchem wir eine wohldefinierte globale Clifford-Multiplikation definieren können. Die Idee ist, einen Modell-Vektorraum $V$ zu wählen, auf diesem alle relevanten Modell-Clifford-Strukturen zu implementieren, und dann nach Wahl geeigneter $G$-Prinzipalbündel $P$ und linearer Darstellung $\lambda:G\to\GL(V)$ der entsprechenden Strukturgruppe $G$ ein assoziiertes Vektorbündel 
$$\Sigma M= P\times_\lambda V$$
bzgl. des Modell-Vektorraum anzusetzen. Dabei soll durch die Wahl des $G$-Prinzipalbündel und der Darstellung sichergestellt werden, dass die Clifford-Struktur des Modell-Vektorraums sich natürlich auf das assoziierte Vektorbündel globalisiert und anschließend die Definition einer globalen Clifford-Multiplikation 
$$c:T^*M\otimes \Sigma M\to \Sigma M$$
zulässt. Dieses assoziierte Vektorbündel wird das \emph{Spinorbündel} $\Sigma M$ sein.

Anschließend müssen wir prüfen, ob wir den angekündigten Spinorzusammenhang $\nabla^\Sigma$ als natürlichen Zusammenhang auf diesem Bündel aus den Riemannschen Daten erhalten und abschließend die Clifford- und Zusammenhangs-Struktur sinnvoll zusammentragen werden kann in der Definition eines natürlichen Differentialoperator 
$$D=c\circ \nabla^\Sigma,$$ dem angekündigten \emph{Dirac-Operator}.

\subsection{Modellhafte Clifford-Strukturen über $(\R^n,\SKP)$ und der Spinorraum $\Sigma_n$}

Wir wollen nun den Modell-Vektorraum mit Clifford-Struktur ansetzen. Dazu wählen wir uns ein möglichst einfaches Setting, welches sich gut in Beziehung setzen lässt zu den Tangentialräumen der Riemannschen Mannigfaltigkeit und deren Clifford-Strukturen.


\begin{DEF}{}{Modellhafte Clifford-Algebra}
Für das euklidische Skalarprodukt auf \(\mathbb{R}^n\) schreiben wir
\[
\Cl_n:=\Cl(\mathbb{R}^n,\SKP).
\]

Im Falle der Komplexifizierung $\C l_n:= \Cl_n\otimes_{\R}\C$ kann bzgl. der Standardbasis $(e_1,…,e_{2m})$ von $V=\R^n$ aus der korrespondierenden $\C$-Basis $(z_1,\cdots,z_m,\bar{z}_1,\cdots,\bar{z}_m)$ von $V\otimes \C$ gegeben durch
\[
z_j:=\frac12(e_{2j-1}-i e_{2j}),\qquad
\bar z_j:=\frac12(e_{2j-1}+i e_{2j})\in \C l_n\qquad  j=1,\dots,m
\] 
eine $\C$-Basis von $\C l_n$ konstruiert werden durch die Produkte
$$ z(\underbrace{(j_1,…,j_k)}_{J_k}) \bar{z}(\underbrace{(i_1,…,i_l)}_{I_\ell})=z_{j_1}\cdots z_{j_k}\,\bar z_{i_1}\cdots \bar z_{i_l},
\qquad k,l=0,\dots,m,\  J_k,I_\ell\subset \{1,\cdots,m\}.$$
\end{DEF}




\begin{DEF}{}{Spinorraum}
Wir definieren bzgl. der angegebenen Basis von $\C l_n$ einen komplexen Untervektorraum von $\C l_n$ durch
\[
\Sigma_n:=
\operatorname{span}_{\C}\bigl\{
z(J_k)\underbrace{(\bar z_1\cdots \bar z_m)}_{\omega}\,\big|\,
k=0,\dots,m,\ J_k \subset \{1,\cdots,m\}
\bigr\}
\subseteq \C l_n.
\]
und nennen diesen den \emph{Spinorraum} in Dimension $n=2m$ und dessen Elemente \emph{Spinoren}.
\end{DEF}


\begin{REM}{}{Dimensionsannahme $n$ gerade}
Hier wird klar, das wir die Dimension der orientierten Riemannschen Mannigfaltigkeit gerade gewählt haben, um die Basis von $\C l_n$ so wählen zu können. Für den ungerade Fall läuft die ganze Konstruktion des Spinorbündels und des Dirac-Operators strukturell analog ab unter Anpassung mancher Schritte, mit welchen wir uns im Rahmen dieser Einführung nicht rumschlagen wollen.    
\end{REM}



\begin{PROP}{}{Spinorraum $\Sigma_n$ ist ein Clifford-Modul über $\C l_n$}
Der Spinorraum $\Sigma_n\subset \C l_n$ ist ein Clifford-Modul über $\C l_n$ mit 
\emph{Modell-Clifford-Multiplikation}
\[
c_0:\R^n\times \Sigma_n\longrightarrow \Sigma_n,
\qquad
(v,\psi)\longmapsto c_0(v)\psi:=v \cdot\psi.
\]
\end{PROP}

\begin{proof}
    Man zeigt zunächst folgende hilfreiche Identitäten für die Basisvektoren von $\Sigma_n$ mit Clifford-Relation:
    \begin{enumerate}
    \item \emph{Case} $l \notin J_k$: 
    $$
    \begin{aligned}
      e_{2 l} \cdot z\left(J_k\right)\omega &=i(-1)^\nu z\left(j_1, \ldots, j_\nu, l, j_{\nu+1}, \ldots, j_k\right)\omega \\
      e_{2 l-1} \cdot z\left(J_k\right)\omega & = (-1)^\nu z\left(j_1, \ldots, j_\nu, l, j_{\nu+1}, \ldots, j_k\right)\omega
    \end{aligned}
    $$
    \item \emph{Case} $l \in J_k$ (wähle $\nu$ sodass $\cdots,j_\nu ,l,j_{\nu+1}\cdots$ in $J_k$):  
    $$
    \begin{aligned}
      e_{2 l} \cdot z\left(j_1, \ldots, j_\nu, l, j_{\nu+1}, \ldots, j_k\right)\omega&=(-1)^\nu i z\left(j_1, \ldots, j_\nu,\hat{l}, j_{\nu+1}, \ldots, j_k\right)\omega\\
      e_{2 l-1} \cdot z\left(j_1, \ldots, j_\nu, l, j_{\nu+1}, \ldots, j_k\right)\omega&=(-1)^{\nu+1} z\left(j_1, \ldots, j_\nu,\hat{l}, j_{\nu+1}, \ldots, j_k\right)\omega .
    \end{aligned}
    $$
    \end{enumerate}
    Exemplarisch für den Fall $l\notin J_k$ zeigt man diese Relationen wie folgt:
    $$
    \begin{aligned}
    e_{2 l} \cdot z\left(J_k\right)\omega & =e_{2 l} \cdot z_{j_1} \cdot \ldots \cdot z_{j_k} \cdot \bar{z}_1 \cdot \ldots \cdot \bar{z}_m \\
    & =(-1)^{k+(l-1)} z_{j_1} \cdot \ldots \cdot z_{j_k} \cdot \bar{z}_1 \cdot \ldots \cdot \bar{z}_{l-1} \cdot e_{2 l} \cdot \bar{z}_l \cdot \bar{z}_{l+1} \cdot \ldots \cdots \bar{z}_m
    \end{aligned}
    $$
    Da
    $$
    \begin{aligned}
    e_{2 l} \cdot \bar{z}_l & =\frac{1}{2} e_{2 l} \cdot\left(e_{2 l-1}+i e_{2 l}\right) \\
    & =\frac{1}{2}\left(e_{2 l} \cdot e_{2 l-1}-i\right) \\
    & =\frac{1}{2}\left(-e_{2 l-1} \cdot e_{2 l}+i e_{2 l-1} \cdot e_{2 l-1}\right) \\
    & =i e_{2 l-1} \cdot \frac{1}{2}\left(e_{2 l-1}+i e_{2 l}\right) \\
    & =i e_{2 l-1} \cdot \bar{z}_l
    \end{aligned}
    $$
    folgt
    $$
    \begin{aligned}
    e_{2 l} \cdot z\left(J_k\right)\omega & =(-1)^{k+l-1} i z_{j_1} \cdot \ldots \cdot z_{j_k} \cdot \bar{z}_1 \cdot \ldots \cdot \bar{z}_{l-1} \cdot e_{2 l-1} \cdot \bar{z}_l \cdot \bar{z}_{l+1} \cdot \ldots \cdot \bar{z}_m \\
    & =(-1)^{k+l-1} i(-1)^{k+l-1} e_{2 l-1} \cdot z_{j_1} \cdot \ldots \cdot z_{j_k} \cdot \bar{z}_1 \cdot \ldots \cdot \bar{z}_m \\
    & =i e_{2 l-1} \cdot z\left(j_1, \ldots, j_k\right)\omega
    \end{aligned}
    $$
    Damit haben wir:
    $$
    \begin{aligned}
    e_{2 l} \cdot z\left(j_1, \ldots, j_k\right)\omega & =\frac{1}{2} e_{2 l} \cdot z\left(j_1, \ldots, j_k\right)\omega+\frac{1}{2} e_{2 l} \cdot z\left(j_1, \ldots, j_k\right)\omega \\
    & = \frac{1}{2} e_{2 l} \cdot z\left(j_1, \ldots, j_k\right)\omega+\frac{i}{2} e_{2 l-1} \cdot z\left(j_1, \ldots, j_k\right)\omega \\
    & =i \underbrace{\frac{1}{2}\left(e_{2 l-1}-i e_{2 l}\right)}_{=z_l} \cdot z\left(j_1, \ldots, j_k\right)\omega \\
    & =i(-1)^\nu z\left(j_1, \ldots, j_\nu, l, j_{\nu+1}, \ldots, j_k\right)\omega
    \end{aligned}
    $$


    Aus den Relationen wird klar, dass $\Sigma_n\subset \C l_n$ unter der Linksmultiplikation mit Elementen aus $\R^n$ invariant ist. Da $\C l_n$ als komplexe Algebra von $\R^n$
    erzeugt wird, folgt daraus auch die Invarianz von $\Sigma_n$ bzgl. Linksmultiplikation mit Elementen aus $\C l_n\cdot \Sigma_n\subseteq \Sigma_n$.

    So erhält man eine wohldefinierter Algebrenhomomorphismus
    \[
    \rho:\C l_n\longrightarrow \End_{\C}(\Sigma_n),
    \qquad
    \rho(z)(\psi):=z\psi.
    \]
    Folglich ist $\Sigma_n$ ein Clifford-Modul über $\C l_n$ mit entsprechender Modell-Clifford-Multiplikation.
\end{proof}



\bigskip 

Wir wollen noch weitere wichtige Eigenschaften des Spinorraums festhalten:

\begin{REM}{}{Eigenschaften des Spinorraums}
\begin{enumerate}
    \item \emph{Clifford-Wirkung auf dem Spinorraum ist vollständig} \parencite[90]{BaerSpin}:
    Die durch Linksmultiplikation induzierte Abbildung
    \[
    \rho:\C l_n\longrightarrow \End_{\C}(\Sigma_n),\qquad
    \rho(X)(\psi):=X\cdot \psi
    \]
    ist sogar ein Isomorphismus komplexer Algebren.
    \item \emph{Chiralitätszerlegung des Spinorraums}:
    Wir definieren für den Spinorraum $\Sigma_n$ die komplexen Untervektorräume
    \[\Sigma_n^{+}:=\operatorname{span}_{\C}\{z(J_k)\omega\mid k=0,\dots,m,\ k \text{ gerade}\},\]
    und
    \[\Sigma_n^{-}:=\operatorname{span}_{\C}\{z(J_k)\omega\mid k=0,\dots,m,\ k \text{ ungerade}\}.\]
    Dann heißt
    \[\Sigma_n=\Sigma_n^{+}\oplus \Sigma_n^{-}\]
    die \emph{Chiralitätszerlegung des Spinorraums}: Die Elemente von $\Sigma_n^\pm$
    heißen Spinoren \emph{positiver} bzw. \emph{negativer Chiralität}.
    \item \emph{Paritätsverhalten unter Clifford-Wirkung}:
    Das Hinzufügen bzw. Entfernen eines $z_l$-Faktors in $z(J_k)$ in den Identitäten der Linkswirkung auf $\Sigma_n$ durch Standardbasisvektoren aus $\R^n$ im Nachweis der Clifford-Moduln-Struktur von $\Sigma_n$ zeigen, dass für die Clifford-Multiplikation auf $\Sigma_n$ bzgl. $\R^n$ Chiralitäten gewechselt werden:
    \[
    \R^n\cdot \Sigma_n^+ \subset \Sigma_n^-,
    \qquad
    \R^n\cdot \Sigma_n^- \subset \Sigma_n^+.
    \]
    Dagegen gilt für den geraden Teil der Clifford-Algebra:
    \[
    \C l_n^0\cdot \Sigma_n^+ \subset \Sigma_n^+,
    \qquad
    \C l_n^0\cdot \Sigma_n^- \subset \Sigma_n^-.
    \]
    \item \emph{Irreduzibilität des Spinorraums}:
    Der \(\C l_n\)-Modul \(\Sigma_n\) ist irreduzibel.

    \begin{proof}
    Sei \(W\subseteq \Sigma_n\) ein nichtverschwindender \(\C l_n\)-Untermodul
    und
    \[
    \psi=\sum_J a_J z(J)\omega \in W\setminus\{0\}.
    \]
    Wähle \(\tilde J\) mit \(a_{\tilde J}\neq 0\) und \(|\tilde J|\) maximal.
    Dann folgt aus den Relationen der \(z_i,\bar z_i\), dass
    \[
    \bar z(\tilde J)\, z(J)\omega = 0 \quad \text{für } J\neq \tilde J
    \]
    und
    \[
    \bar z(\tilde J)\, z(\tilde J)\omega = \pm \omega.
    \]
    Also
    \[
    \bar z(\tilde J)\psi = \pm a_{\tilde J}\omega \neq 0.
    \]
    Damit ist $\omega \in W$. Durch Linksmultiplikation mit den Produkten \(z(J_k)\) folgt dann \[ z(J_k)\omega\in W. \qquad \] Damit enthält \(W\) eine Basis von \(\Sigma_n\), also \(W=\Sigma_n\).
    \end{proof}
\end{enumerate}
\end{REM}



\subsection{Zusammenhang zwischen $\C l_n$ und Spin Gruppe}


Wir stellen uns jetzt die Frage, was die natürliche Strukturgruppe $G$ sein könnte, welche eine lineare Darstellung im Spinorraum $\Sigma_n$ zulässt. Wir werden sehen, dass hier die \emph{Spin-Gruppe} ihren natürlichen Auftritt hat. 

\begin{DEF}{}{Pin- und Spin-Gruppe}
\begin{enumerate}
\item Die \emph{Pin-Gruppe} ist definiert als die von den Einheitsvektoren erzeugte Untergruppe
\[
\Pin(n):=\langle S^{n-1}\rangle = \left\{v_1 \cdot \ldots \cdot v_m \in \mathrm{Cl}_n \mid v_j \in S^{n-1} \subset \mathbb{R}^n, m \in \mathbb{N}_0\right\} \subset \Cl_n^\times.
\]

\item Die \emph{Spin-Gruppe} ist der gerade Teil von der Pin Gruppe:
$$
\begin{aligned}
\operatorname{Spin}(n) & :=\operatorname{Pin}(n) \cap \mathrm{Cl}_n^0 \\
& =\left\{v_1 \cdot \ldots \cdot v_m \in \mathrm{Cl}_n \mid v_j \in S^{n-1}, m \in 2 \mathbb{N}_0\right\} .
\end{aligned}
$$
\end{enumerate}
\end{DEF}

\bigskip

Wir halten auch direkt wichtige Eigenschaften fest:


\begin{REM}{}{Eigenschaften der Spin Gruppe}
  \begin{enumerate}
    \item \emph{Grundlegendes}: $\Spin(n)$ ist kompakt, zusammenhängend für $n\geq 2$ und einfach zusammenhängend für $n\geq 3$.
    \item \emph{Lie-Untergruppe}: $\Spin(n)$ ist eine abgeschlossene Lie-Untergruppe von $\Cl_n^\times$ und $\Pin(n)$.
    \item \emph{Zweifache Überlagerung}: Man definiert durch  
    \[
    \Lambda:\Pin(n)\to \mathrm{O}(n);\ u\mapsto [x\mapsto \Lambda(u)(x):=\alpha(u)\,x\,u^{-1}]
    \]
    ein surjektiver Gruppenhomomorphismus mit $\ker\Lambda=\{\pm 1\}$, wobei $\alpha:\Cl_n\to\Cl_n$ die Gradinvolution ist, d.h.
    \[
    \alpha(v)=-v \quad \text{für } v\in \R^n,
    \qquad
    \alpha(x)=(-1)^k x \text{ für } x\in \Cl_n^{(k)}.
    \]
    Einschränkung von $\Lambda$ auf die Spin Gruppe $\Spin(n)\subset\Pin(n)$ liefert einen surjektiven Liegruppenhomomorphismus
    \[
    \lambda:=\Lambda|_{\Spin(n)}:\Spin(n)\to \SO(n),
    \qquad
    \lambda(u)(x)=u x u^{-1},
    \]
    auch mit $\ker\lambda=\{\pm 1\}$, die kanonische zweifache Überlagerung von $\SO(n)$ durch $\Spin(n)$.
    \item \emph{Lie Algebra}: Die Lie Algebra der Spin Gruppe $\mathfrak{spin}(n)$ ist gegeben durch
    $$\mathfrak{spin}(n)=\operatorname{span}_{\mathbb{R}}\left\{e_i e_j \mid i<j\right\} \subset \mathrm{Cl}_n$$
    und hat Dimension $\dim \mathfrak{spin}(n)=\frac{n(n-1)}{2},$
    was durch Kurven 
    $$c: \mathbb{R} \rightarrow \operatorname{Spin}(n);\ t \mapsto\left(\cos (t) e_i+\sin (t) e_j\right) \cdot\left(-e_i\right) \qquad c(0)=1, \quad c^{\prime}(0)=e_i e_j$$
    geprüft werden kann.
    \item \emph{Isomorphe Lie Algebren}: Es gilt für das Lie Differential der zweiblättrigen Überlagerung
    $$d\lambda:\mathfrak{spin}(n)\overset{\cong}{\longrightarrow}\mathfrak{so}(n)$$
    
  \end{enumerate}
\end{REM}

\bigskip

Da $\Spin(n)\subset \Cl_0\subset\C l_n$, erhalten wir eine natürliche Darstellung der Spin Gruppe im Spinorraum durch Einschränkung:

\begin{DEF}{}{Spinor- und Halbspin-Darstellungen}
Die Einschränkung der Clifford-Darstellung 
$$\rho: \C l_n\to \End (\Sigma_n)$$
auf die Spin Gruppe induzierte die \emph{Spinor-Darstellung}
\[
\sigma_n:=\rho|_{\Spin(n)}:\Spin(n)\to \GL(\Sigma_n)
\]
als lineare Darstellung.

In gerader Dimension induziert die chirale Zerlegung
\[
\Sigma_n=\Sigma_n^+\oplus \Sigma_n^-
\]
die Darstellungen
\[
\sigma_n^\pm:\Spin(n)\to \GL(\Sigma_n^\pm),
\]
die \emph{positive} bzw. \emph{negative Halbspin-Darstellung} heißen. In gerader Dimension ist die Spinor-Darstellung also reduzibel.
\end{DEF}







\subsection{Zusammenhang zwischen geometrischer und modellhafter Clifford-Struktur}


Wir klären jetzt den natürlichen Zusammenhang zwischen der geometrischen und Modell-Clifford-Strukturen.

\begin{LEM}{}{Modellbeschreibung von $\C l_n(T_xM,g_x)$ durch $\C l_n$} Isometrien zwischen $(\R^n,\SKP)$ und $(T_xM,g_x)$ induzieren Algebrenisomorphismen $\C l(u): \C l_n \overset{\cong}{\to} \C l(T_xM,g_x)$.
\end{LEM}


\begin{proof}
    
Jeder orientierte orthonormale Rahmen $u=(u_1,\cdots,u_n)$ über $x$ liefert für $u_i:=u(e_i)$ eine Isometrie 
$$u\in P_x^{\SO}(M):=\left\{u:\R^n\to T_xM\;\middle|\;u \text{ ist orientierungserhaltende Isometrie}\right\}$$ 
durch
\[
u:\bigl(\R^n,\SKP\bigr)\xrightarrow{\;\cong\;}(T_xM,g_x),
\qquad
(a^1,\dots,a^n)\longmapsto \sum_{i=1}^n a^i u_i.
\]
Nach der universellen Eigenschaft der Clifford-Algebra induziert \(u\) einen
eindeutigen Algebrenhomomorphismus
\[
\C l(u):\C l_n\longrightarrow \C l(T_xM,g_x),
\]
so dass das Diagramm
\[
\begin{tikzcd}
\R^n \arrow[r, "\iota_{\mathrm{euc}}"] \arrow[d, "u"'] 
& \C l_n \arrow[d, dashed, "\C l(u)"] \\
T_xM \arrow[r, "\iota_x"'] 
& \C l(T_xM,g_x)
\end{tikzcd}
\]
kommutiert.

Da \(u\) eine Isometrie ist, ist auch \(u^{-1}\) eine Isometrie und induziert
analog einen Homomorphismus
\[
\C l(u^{-1}):\C l(T_xM,g_x)\to \C l_n.
\]
Wegen der Eindeutigkeit in der universellen Eigenschaft sind
\[
\C l(u^{-1})\circ \C l(u)=\id_{\C l_n},
\qquad
\C l(u)\circ \C l(u^{-1})=\id_{\C l(T_xM,g_x)},
\]
also ist \(\C l(u^{-1})\) ein Algebraisomorphismus
$$\C l(u^{-1}): \C l(T_xM,g_x)\overset{\cong}{\to} \C l_n$$
welcher aber jeweils von der Wahl eines Rahmens $u$ abhängt und daher zunächst nicht kanonisch gesetzt werden kann. In einem orthonormalen Rahmen werden die punktweisen Clifford-Algebren also alle bis auf Isomorphie durch dasselbe Modell \(\C l_n\) beschrieben. 

\end{proof}


\bigskip

Aus $P_x^\SO (M)$ als Faser lässt sich auch ein natürliches $\SO(n)$-Prinzipalbündel bauen, das sogenannte $\SO(n)$-Rahmenbündel.


\begin{PROP}{}{$\SO(n)$-Rahmenbündel ist $\SO(n)$-Prinzipalbündel}
Sei $(M,g)$ eine orientierte Riemannsche Mannigfaltigkeit. Für \(x\in M\) setzen wir
\[
P_x^{\SO}(M)
:=
\left\{
h:\R^n\to T_xM
\;\middle|\;
h \text{ ist orientierungserhaltende Isometrie}
\right\}.
\]
Ferner definieren wir
\[
P^{\SO}(M):=\bigsqcup_{x\in M} P_x^{\SO}(M)
\]
und die Projektion
\[
\pi:P^{\SO}(M)\to M,\qquad \pi(h)=x
\quad\text{für }h\in P_x^{\SO}(M).
\]
Auf \(P_x^{\SO}(M)\) wirkt \(\SO(n)\) von rechts durch
\[
P_x^{\SO}(M)\times \SO(n)\to P_x^{\SO}(M),
\qquad
(h,A)\mapsto h\circ A.
\]
Dann ist
\[
\pi:P^{\SO}(M)\to M
\]
ein \(\SO(n)\)-Prinzipalbündel über $M$, welches wir als $\SO(n)$-\emph{Rahmenbündel von $(M,g)$} bezeichnen. Über jedem \(x\in M\) ist die Faser \(P_x^{\SO}(M)\) ein freier transitiver \(\SO(n)\)-Raum.
\end{PROP}

\begin{proof}

Die Rechtswirkung
\[
R_A:P^{\SO}(M)\to P^{\SO}(M),\qquad h\mapsto h\circ A,
\qquad A\in \SO(n),
\]
ist wohldefiniert, denn ist \(h\in P_x^{\SO}(M)\), so ist auch
\[
h\circ A:\R^n\to T_xM
\]
wieder eine orientierungserhaltende Isometrie. Insbesondere gilt
\[
\pi(h\circ A)=\pi(h),
\]
also ist die Wirkung fasererhaltend.

Außerdem wirkt \(\SO(n)\) auf jeder Faser einfach transitiv: Sind
\(h_1,h_2\in P_x^{\SO}(M)\), so ist
\[
A:=h_1^{-1}\circ h_2\in \SO(n)
\]
und
\[
h_1\circ A=h_2.
\]
Damit ist die Wirkung transitiv. Ist umgekehrt \(h\circ A=h\), so folgt wegen
der Invertierbarkeit von \(h\),
\[
A=\id_{\R^n},
\]
also ist die Wirkung frei.

Es bleiben die lokalen \(\SO(n)\)-äquivarianten Trivialisierungen. Sei
\(U\subset M\) offen und
\[
s=(s_1,\dots,s_n)
\]
ein lokaler orientierter orthonormaler Rahmen auf \(U\). Dann definiert
\[
s(x):\R^n\to T_xM,\qquad e_i\mapsto s_i(x),
\]
einen Schnitt \(s:U\to P^{\SO}(M)\). Wegen der einfachen Transitivität der
Wirkung auf jeder Faser besitzt jedes \(h\in \pi^{-1}(U)\) eindeutig die Form
\[
h=s(\pi(h))\circ A
\qquad\text{mit }A\in \SO(n).
\]
Somit erhält man eine lokale Trivialisierung
\[
\Phi_U:\pi^{-1}(U)\to U\times \SO(n),\qquad
h\mapsto \bigl(\pi(h),A\bigr),
\]
mit Inverser
\[
\Phi_U^{-1}(x,A)=s(x)\circ A.
\]
Aus
\[
\Phi_U(h\circ B)=(x,AB)=\Phi_U(h)\cdot B
\qquad (B\in \SO(n))
\]
folgt schließlich die \(\SO(n)\)-Äquivarianz.

Damit sind alle Bedingungen eines \(\SO(n)\)-Prinzipalbündels erfüllt.
\end{proof}





\subsection{Spin-Struktur, Spinorbündel und globale Clifford-Multiplikation}


Wir können jetzt nach unserem Ansatz die eingeführten Strukturen zusammentragen in einem assoziierten Vektorbündel. Da unsere Strukturgruppe offensichtlich die Spin-Gruppe sein wird, müssen wir dabei darauf achten, dass zusätzlich ein verträglicher Zusammenhang 
$$\Theta: P^\Spin(M)\to P^\SO(M)$$
zwischen möglichen $\Spin(n)$-Prinzipalbündeln über $M$ und dem $\SO(n)$-Rahmenbündel sichergestellt wird, mittels welchem wir in der Definition der globalen Clifford-Multiplikation via $\Theta(q)\in P^\SO(M)$ durch 
$$\Theta(q)^{-1}:T_xM\to \R^n$$
die faserweise Clifford-Multiplikation bzgl. Tangentialvektoren $v\in T_xM$ übersetzen können in unsere Modell-Clifford-Multiplikation $c_0$ bzgl. Vektoren $\Theta(q)^{-1}(v)\in\R^n$.

\begin{DEF}{}{Riemannsche Spin-Mannigfaltigkeit und Spinorbündel}
Sei \(M\) eine orientierte Riemannsche Mannigfaltigkeit. 

\begin{enumerate}

\item \emph{Spin-Struktur}: Eine \emph{Spin-Struktur}
auf \(M\) ist ein Paar
\[
\bigl(P^{\Spin}(M),\Theta\bigr),
\]
bestehend aus
\begin{enumerate}
    \item einem \(\Spin(n)\)-Prinzipalbündel
    \[
    \pi_{\Spin}:P^{\Spin}(M)\to M
    \]
    über \(M\),
    \item einer zweiblättrigen Überlagerung
    \[
    \Theta:P^{\Spin}(M)\to P^{\SO}(M),
    \]
    so dass das folgende Diagramm kommutiert:
    \[
    \begin{tikzcd}[column sep=large, row sep=large]
    P^{\Spin}(M)\times \Spin(n) \arrow[r,"R_{\Spin(n)}"] \arrow[d,"\Theta\times \lambda"'] &
    P^{\Spin}(M) \arrow[rdd,"\pi_{\Spin}"] \arrow[d,"\Theta"] & \\
    P^{\SO}(M)\times \SO(n) \arrow[r,"R_{\SO(n)}"] &
    P^{\SO}(M) \arrow[rd,"\pi_{\SO}"'] & \\
    & & M
    \end{tikzcd}
    \]
    Dabei gilt insbesondere die Äquivarianzbedingung
    $$\Theta(Hg)=\Theta(H)\lambda(g)$$
    Hier bezeichnet
    \[
    \lambda:\Spin(n)\to \SO(n)
    \]
    die zweiblättrige Überlagerung von \(\SO(n)\). Die horizontalen Pfeile sind
    jeweils die Rechtswirkungen der Strukturgruppen auf den entsprechenden
    Prinzipalbündeln.
\end{enumerate}

Eine orientierte Riemannsche Mannigfaltigkeit, die mit einer Spin-Struktur
ausgestattet ist, heißt \emph{Riemannsche Spin-Mannigfaltigkeit}, oder kurz \emph{Spin-Mannigfaltigkeit}. Eine orientierte Riemannsche Mannigfaltigkeit heißt \emph{spinnbar}, wenn sie eine Spin-Struktur zulässt.

\item \emph{Spinorbündel}: Sei $\sigma_n:\Spin(n)\to \GL(\Sigma_n)$ die Spinor-Darstellung. Das \emph{Spinorbündel} zu einer Riemannschen Spin-Mannigfaltigkeit mit Spin-Struktur $(P^\Spin(M),\Theta)$ ist das assoziierte Vektorbündel
\[
\Sigma M:=P^{\Spin}(M)\times_{\sigma_n}\Sigma_n
\]
das \emph{Spinorbündel} von \(M\) bezüglich der gegebenen Spin-Struktur mit Äquivalenzklassen 
$$[H,\phi]\in\Sigma M$$
als Elemente, wobei $H\in P^{\Spin}(M)$ ein \emph{Spin-Rahmen} und $\phi\in \Sigma_n$ ein \emph{Spinor} ist. Schnitte von \(\Sigma M\) heißen \emph{Spinorfelder} auf \(M\).

Für den Fall \(n\) gerade und $\sigma_n^\pm:\Spin(n)\to \GL(\Sigma_n^\pm)$
die positive bzw. negative Halbspin-Darstellung heißen die assoziierten
Vektorbündel
\[
\Sigma^\pm M:=P^{\Spin}(M)\times_{\sigma_n^\pm}\Sigma_n^\pm
\]
das \emph{positive} bzw. \emph{negative Spinorbündel} von \(M\) bezüglich der
gegebenen Spin-Struktur.

\end{enumerate}

\end{DEF}

\begin{THEO}{}{Existenz der globalen Clifford-Multiplikation}
Sei $(M,g)$ eine orientierte Riemannsche Spin-Mannigfaltigkeit mit Spin-Struktur und Spinorbündel $\Sigma M$. 

Dann existiert bzgl. des Faserproduktes 
$$TM\times_M \Sigma M= \coprod_{x \in M}\left(T_x M \times \Sigma_x M\right)$$
eine glatte Bündelabbildung
\[
c:TM\times_M \Sigma M\longrightarrow \Sigma M,
\]
deren Einschränkung auf jede Faser $x\in M$ eine bilineare Abbildung
\[
c_x:T_xM\times \Sigma_xM\longrightarrow \Sigma_xM
\]
liefert, so dass für alle $v,w\in T_xM$ gilt:
\[
c_x(v)c_x(w)+c_x(w)c_x(v)
=
-2g_x(v,w)\id_{\Sigma_xM}.
\]

Dabei wird $c_x$ durch die Modell-Clifford-Multiplikation definiert
\[
c_x(v)[q,\psi]
=
[q,\,c_0(u_q^{-1}(v))\psi],
\]
wobei $q\in P_x^{\Spin}$ und
\[
u_q:=\Theta(q):\R^n\to T_xM
\]
der durch $q$ induzierte orthonormale Rahmen ist.
\end{THEO}

\begin{proof}
Die Äquivarianz der Spin-Struktur,
\[
\Theta(qg)=\Theta(q)\cdot \lambda(g)
\qquad g\in \Spin(n),
\]
impliziert
\[
u_{qg}=u_q\circ \lambda(g)
\qquad\text{und daher}\qquad
u_{qg}^{-1}=\lambda(g)^{-1}\circ u_q^{-1}.
\]

Des Weiteren gilt für $\psi\in \Sigma_n$ wegen der Modulstruktur und der Definition von $\rho$:
\[
\begin{aligned}
\rho(g)\,c_0(v)\,\rho(g)^{-1}(\psi)
&=
g\cdot \bigl(v\cdot (g^{-1}\cdot \psi)\bigr) \\
&=
(gvg^{-1})\cdot \psi \\
&=
c_0(\lambda(g)v)\psi.
\end{aligned}
\]

\begin{enumerate}
    \item \emph{Definition und Wohldefiniertheit.}
    Für $x\in M$, $v\in T_xM$ und $[q,\psi]\in \Sigma_xM$ definieren wir
    \[
    c_x(v)[q,\psi]
    :=
    [q,\,c_0(u_q^{-1}(v))\psi].
    \]
    Es ist zu zeigen, dass diese Definition unabhängig von der Wahl des Repräsentanten
    von $[q,\psi]$ ist.

    Sei dazu $g\in \Spin(n)$. Dann gilt
    \[
    [q,\psi]=[qg,\rho(g^{-1})\psi].
    \]
    Also
    \[
    \begin{aligned}
    c_x(v)[qg,\rho(g^{-1})\psi]
    &=
    [qg,\,c_0(u_{qg}^{-1}(v))\,\rho(g^{-1})\psi] \\
    &=
    [qg,\,c_0(\lambda(g)^{-1}u_q^{-1}(v))\,\rho(g^{-1})\psi].
    \end{aligned}
    \]
    Anwendung der $\Spin(n)$-Äquivarianz der Modell-Clifford-Multiplikation bzgl. $u_q^{-1}(v)$ gibt
    \[
    \begin{aligned}
    c_x(v)[qg,\rho(g^{-1})\psi]
    &=
    [qg,\rho(g^{-1})\bigl(c_0(u_q^{-1}(v))\psi\bigr)] \\
    &=
    [q,\,c_0(u_q^{-1}(v))\psi] \\
    &=
    c_x(v)[q,\psi].
    \end{aligned}
    \]
    Also ist $c_x$ wohldefiniert.

    \item \emph{Bilinearität und Glattheit.}
    Die Bilinearität der faserweisen Abbildungen
    \[
    c_x:T_xM\times \Sigma_xM\to \Sigma_xM
    \]
    folgt unmittelbar aus der Bilinearität der Modell-Clifford-Multiplikation.

    Zur Prüfung der Glattheit wählen wir auf einer offenen Menge $U\subset M$ einen glatten
    lokalen Schnitt
    \[
    s:U\to P^{\Spin}(M).
    \]
    Da
    \[
    TM\cong P^{\Spin}\times_\lambda \R^n,
    \qquad
    \Sigma M\cong P^{\Spin}\times_\rho \Sigma_n,
    \]
    induziert $s$ Bündelisomorphismen
    \[
    \Phi_s^{TM}:U\times \R^n\to TM|_U,
    \qquad
    (x,\xi)\mapsto [s(x),\xi],
    \]
    \[
    \Phi_s^{\Sigma}:U\times \Sigma_n\to \Sigma M|_U,
    \qquad
    (x,\psi)\mapsto [s(x),\psi].
    \]
    Damit erhält man eine lokale Darstellung
    \[
    \widetilde c_s
    :=
    (\Phi_s^\Sigma)^{-1}\circ c\circ (\Phi_s^{TM}\times_U \Phi_s^\Sigma)
    :
    U\times \R^n\times \Sigma_n\to U\times \Sigma_n.
    \]
    Für $(x,\xi,\psi)\in U\times \R^n\times \Sigma_n$ gilt wegen der Definition von $c$:
    \[
    \begin{aligned}
    \widetilde c_s(x,\xi,\psi)
    &=
    (\Phi_s^\Sigma)^{-1}\bigl(c([s(x),\xi],[s(x),\psi])\bigr)\\
    &=
    (\Phi_s^\Sigma)^{-1}\bigl([s(x),c_0(\xi)\psi]\bigr)\\
    &=
    (x,c_0(\xi)\psi).
    \end{aligned}
    \]
    In diesen lokalen Trivialisierungen wird $c$ also durch
    \[
    (x,\xi,\psi)\mapsto (x,c_0(\xi)\psi)
    \]
    beschrieben. Da $c_0$ bilinear und damit glatt ist, ist auch die lokale Darstellung
    glatt. Ist die Clifford-Multiplikation eine glatte Bündelabbildung.
    \item \emph{Clifford-Relation.}
    Für $x\in M$, $v,w\in T_xM$ und $[q,\psi]\in \Sigma_xM$ gilt
    \[
    \begin{aligned}
    c_x(v)\bigl(c_x(w)[q,\psi]\bigr)
    &=
    c_x(v)[q,\,c_0(u_q^{-1}(w))\psi] \\
    &=
    [q,\,c_0(u_q^{-1}(v))\bigl(c_0(u_q^{-1}(w))\psi\bigr)].
    \end{aligned}
    \]
    Daher
    \[
    \begin{aligned}
    &\bigl(c_x(v)c_x(w)+c_x(w)c_x(v)\bigr)[q,\psi] \\
    &=
    [q,\,
    \bigl(c_0(u_q^{-1}(v))c_0(u_q^{-1}(w))
    +c_0(u_q^{-1}(w))c_0(u_q^{-1}(v))\bigr)\psi].
    \end{aligned}
    \]
    Aus der Clifford-Relation des Modellraums folgt
    \[
    \begin{aligned}
    \bigl(c_x(v)c_x(w)+c_x(w)c_x(v)\bigr)[q,\psi] =-2\langle u_q^{-1}(v),u_q^{-1}(w)\rangle\,[q,\psi].
    \end{aligned}
    \]
    Da $u_q:\R^n\to T_xM$ eine lineare Isometrie ist, gilt
    \[
    \langle u_q^{-1}(v),u_q^{-1}(w)\rangle=g_x(v,w).
    \]
    Also
    \[
    c_x(v)c_x(w)+c_x(w)c_x(v)
    =
    -2g_x(v,w)\id_{\Sigma_xM}.
    \]

\end{enumerate}
\end{proof}

\begin{REM}{}{Tensorielle Form der Clifford-Multiplikation}
Da die globale Clifford-Multiplikation
\[
c:TM\times_M \Sigma M\to \Sigma M
\]
faserweise bilinear ist, induziert sie eindeutig eine Bündelabbildung
\[
c:TM\otimes \Sigma M\to \Sigma M.
\]
Mittels des musikalischen Isomorphismus $\sharp:T^*M\to TM$ kann sie ebenso als
Bündelabbildung
\[
c:T^*M\otimes \Sigma M\to \Sigma M
\]
aufgefasst werden.
\end{REM}


\subsection{Existenz und Äquivalenz von Spin-Strukturen}


Wir wollen jetzt kurz auf die Frage eingehen, wann überhaupt Spin-Strukturen auf einer Riemannschen Mannigfaltigkeit existieren, und sofern welche existieren, welche äquivalent sind bzw. wie viele Klassen von nicht-äquivalenten Spin-Strukturen es gibt. Die Details und Beispiele findet man u.a. in \parencite[p.~85-93]{LawsonMichelsohn2016}. 


\begin{DEF}{}{Äquivalente Spin-Strukturen}
Seien
\[
\bigl(P_1^{\Spin}(M),\Theta_1\bigr)
\qquad\text{und}\qquad
\bigl(P_2^{\Spin}(M),\Theta_2\bigr)
\]
zwei Spin-Strukturen auf \(M\). Diese heißen \emph{äquivalent}, wenn es einen
\(\Spin(n)\)-äquivarianten Diffeomorphismus
\[
\phi:P_1^{\Spin}(M)\to P_2^{\Spin}(M)
\]
gibt, so dass das folgende Diagramm kommutiert:
\[
\begin{tikzcd}[column sep=large, row sep=large]
P_1^{\Spin}(M)\times \Spin(n) \arrow[r,"\phi\times \id"] \arrow[d,"R_{\Spin(n)}"] &
P_2^{\Spin}(M)\times \Spin(n) \arrow[d,"R_{\Spin(n)}"] \\
P_1^{\Spin}(M) \arrow[r,"\phi"] \arrow[dr,"\Theta_1"'] &
P_2^{\Spin}(M) \arrow[d,"\Theta_2"] \\
& P^{\SO}(M)
\end{tikzcd}
\]
Dabei sind die vertikalen Pfeile oben die jeweiligen Rechtswirkungen von
\(\Spin(n)\) auf \(P_1^{\Spin}(M)\) bzw. \(P_2^{\Spin}(M)\).
\end{DEF}

\bigskip

Wir gehen zunächst das einfache Beispiel der Sphäre $S^1$ durch:

\begin{EXA}{}{Spin-Strukturen auf $S^1$}
Für $M=S^1$ haben wir $\mathrm{SO}(1)=\{1\}$ und daher $\operatorname{Spin}(1)=\mathbb{Z}_2$ und $P^{\mathrm{SO}}\left(S^1\right)=S^1$ trivial, da es an jedem Punkt genau einen orientierten ON-Rahmen gibt. Eine Spin-Struktur ist daher ein $\mathbb{Z}_2$-Hauptbündel über $S^1$, also eine zweifache Überlagerung. Es gibt zwei Isomorphietypen:
\begin{enumerate}
    \item \emph{Triviale Spin-Struktur}:
    $$
    P_{\text {triv }}^{\text {Spin }}\left(S^1\right)=S^1 \sqcup S^1=S^1 \times \mathbb{Z}_2 .
    $$
    Sei $\Theta: P^{\operatorname{Spin}}\left(S^1\right)=S^1 \times \mathbb{Z}_2 \rightarrow P^{\mathrm{SO}}\left(S^1\right)=S^1$ die Projektion auf die erste Komponente. Das entsprechende Diagramm kommutiert:
    $$\begin{tikzcd}
    S^1 \times \mathbb{Z}_2 \times \mathbb{Z}_2 
    \arrow[r, "{\mathrm{id} \times \mu}"] 
    \arrow[d, "{\mathrm{pr}_1 \times \lambda}"'] 
    & S^1 \times \mathbb{Z}_2 
    \arrow[d, "{\mathrm{pr}_1}"] 
    \arrow[dr, "{\mathrm{pr}_1}"] \\
    S^1 \times \{1\} 
    \arrow[r, "{\mathrm{pr}_1}"'] 
    & S^1 
    \arrow[r, "\mathrm{id}"'] 
    & S^1
    \end{tikzcd}$$
    
    \item \emph{Nicht-triviale Spin-Struktur}:
    $$P_{\text {non-triv }}^{\text {Spin }}\left(S^1\right)=S^1$$ 
    mit $\Theta: S^1 \rightarrow S^1, z \mapsto z^2$. Die Rechtswirkung von $\operatorname{Spin}(1)=\mathbb{Z}_2$ auf $P_{\text {non-triv }}^{\text {Spin }}\left(S^1\right)$ ist gegeben durch $z \mapsto-z$. Das entsprechende Diagramm kommutiert
    $$\begin{tikzcd}
    S^1 \times \mathbb{Z}_2 
    \arrow[r, "{(z,\varepsilon)\mapsto \varepsilon z}"] 
    \arrow[d, "{(z,\varepsilon)\mapsto (z^2,1)}"'] 
    & S^1 
    \arrow[d, "{z\mapsto z^2}"] 
    \arrow[dr, "{z\mapsto z^2}"] \\
    S^1 \times \{1\} 
    \arrow[r, "{\mathrm{pr}_1}"'] 
    & S^1 
    \arrow[r, "\mathrm{id}"'] 
    & S^1
    \end{tikzcd}$$
\end{enumerate}
Beide Spin-Strukturen sind offensichtlich nicht äquivalent.
\end{EXA}

\bigskip 

Das klassische Resultat zur Untersuchung der Existenz- und Äquivalenz-Fragen für Spin-Strukturen ist:

\begin{THEO}{\parencite[p.~86]{LawsonMichelsohn2016}}{Existenz und Äquivalenz von Spin-Strukturen}
    Eine orientierte Riemannsche Mannigfaltigkeit $M$ lässt Spin-Strukturen zu genau dann wenn die zweite Stiefel-Whitney-Klasse $w_2(TM) \in H^2(M ; \mathbb{Z}_2)$ verschwindet, also $w_2(TM)=0$. Dann stehen die Äquivalenzklassen von Spin-Strukturen in bijektiver Beziehung zu den Elementen von $H^1\left(M; \mathbb{Z}_2\right)$.
\end{THEO}



\begin{EXA}{\parencite[p.~85-93]{LawsonMichelsohn2016}}{Klassische Beispiele Riemannscher Spin-Mannigfaltigkeiten} 
Zu den wichtigsten Klassen Riemannscher Spin-Mannigfaltigkeiten gehören:

\begin{enumerate}
    \item \textbf{Der euklidische Raum \(\R^n\).}  
    Da \(T\R^n\) trivial ist, ist auch das orientierte orthonormale Rahmenbündel trivial, und \(\R^n\) besitzt eine Spin-Struktur.

    \item \textbf{Der Torus \(T^n=\R^n/\Z^n\).}  
    Auch hier ist das Tangentialbündel trivial, also ist \(T^n\) spin. Allgemeiner ist jede parallelisierbare orientierte Mannigfaltigkeit spin.

    \item \textbf{Die Sphäre \(S^n\) für \(n\ge 2\).}  
    Die Standardsphäre mit ihrer Riemannschen Metrik ist orientierbar und besitzt eine Spin-Struktur. Sie ist eines der grundlegenden Beispiele kompakter Spin-Mannigfaltigkeiten.

    \item \textbf{Jede orientierte Fläche.}  
    Für eine orientierte zweidimensionale Riemannsche Mannigfaltigkeit \(M\) gilt stets $w_2(M)=0$, also ist jede orientierte Fläche spin. Insbesondere sind \(S^2\), \(T^2\) und höhergenusige geschlossene orientierte Flächen Spin-Mannigfaltigkeiten.

    \item \textbf{Komplex projektive Räume \(\CP^m\).}  
    Diese sind genau dann spin, wenn \(m\) ungerade ist. Insbesondere ist \(\CP^1\cong S^2\) spin, während \(\CP^2\) kein Spin-Beispiel ist.

    \item \textbf{Calabi--Yau-Mannigfaltigkeiten.}  
    Jede glatte Calabi--Yau-Mannigfaltigkeit ist spin. Denn für eine komplexe Mannigfaltigkeit gilt $w_2(M)\equiv c_1(T^{1,0}M)\mod 2$, und bei einer Calabi--Yau-Mannigfaltigkeit ist das kanonische Bündel \(K_M\) trivial, also $c_1(T^{1,0}M)=-c_1(K_M)=0$. Daher folgt $w_2(M)=0$. Wichtige Beispiele sind elliptische Kurven, \(K3\)-Flächen und komplexe Tori.
    
    \item \textbf{Lie-Gruppen mit linksinvarianter Riemannscher Metrik.}  
    Jede zusammenhängende Lie-Gruppe ist parallelisierbar, also im orientierbaren Fall spin. Damit liefern etwa \(\SU(2)\cong S^3\), \(\SO(3)\), \(T^n\) oder allgemein viele kompakte Lie-Gruppen natürliche Beispiele von Spin-Mannigfaltigkeiten.

    \item \textbf{Produkte von Spin-Mannigfaltigkeiten.}  
    Sind \(M\) und \(N\) Riemannsche Spin-Mannigfaltigkeiten, so ist auch das Produkt \(M\times N\) mit der Produktsmetrik wieder spin.
\end{enumerate}

\end{EXA}

\begin{EXA}{\parencite[p.~85-93]{LawsonMichelsohn2016}}{Klassische Beispiele nicht-spinischer Riemannscher Mannigfaltigkeiten}

\begin{enumerate}
    \item \textbf{Der komplex projektive Raum \(\CP^{2k}\).}  
    Für komplex projektive Räume ist $\CP^m$ genau dann spin, wenn $m$ ungerade ist.
    Insbesondere ist also \(\CP^{2k}\) für jedes \(k\ge 1\) nicht spin.

    \item \textbf{Nichtorientierbare Riemannsche Mannigfaltigkeiten.}  
    Jede Spin-Mannigfaltigkeit ist insbesondere orientierbar. Daher kann eine nichtorientierbare Riemannsche Mannigfaltigkeit keine Spin-Struktur tragen. Beispiele sind etwa $\RP^{2}$, das Möbiusband, die Kleinsche Flasche.

    \item \textbf{Reell projektive Räume \(\RP^n\) in vielen Dimensionen.}  
    Zunächst ist \(\RP^n\) genau dann orientierbar, wenn \(n\) ungerade ist. Unter den orientierbaren Fällen ist \(\RP^n\) genau dann spin, wenn $n\equiv 3 \pmod 4$. Somit sind etwa $\RP^5,\ \RP^9,\ \RP^{13},\dots$ orientierbar, aber nicht spin.
\end{enumerate}

\end{EXA}





\begin{EXA}{\parencite[p.~85-93]{LawsonMichelsohn2016}}{Beispiele für eindeutige bzw. mehrere Spin-Strukturen}
\begin{enumerate}
    \item \textbf{Eindeutige Spin-Struktur auf \(S^n\) (\(n\ge 2\)).}  
    Die Sphäre \(S^n\) ist spin und erfüllt $H^1(S^n;\Z_2)=0$. Daher besitzt \(S^n\) genau eine Spin-Struktur.

    \item \textbf{Eindeutige Spin-Struktur auf einfach zusammenhängenden Spin-Mannigfaltigkeiten.}  
    Allgemeiner besitzt jede einfach zusammenhängende Spin-Mannigfaltigkeit genau eine Spin-Struktur, da dann $H^1(M;\Z_2)=0$ gilt. Beispiele sind etwa \(S^n\), \(K3\)-Flächen oder allgemein viele einfach zusammenhängende Calabi--Yau-Mannigfaltigkeiten.

    \item \textbf{Eindeutige Spin-Struktur auf \(\CP^{2m+1}\).}  
    Die komplex projektiven Räume \(\CP^{2m+1}\) sind spin, und wegen $H^1(\CP^{2m+1};\Z_2)=0$ besitzen sie genau eine Spin-Struktur.

    \item \textbf{Zwei Spin-Strukturen auf \(S^1\).}  
    Der Kreis ist spin und erfüllt $H^1(S^1;\Z_2)\cong \Z_2$. Daher besitzt \(S^1\) genau zwei verschiedene Spin-Strukturen.

    \item \textbf{\(2^n\) Spin-Strukturen auf dem Torus \(T^n\).}  
    Der \(n\)-dimensionale Torus ist spin und hat $H^1(T^n;\Z_2)\cong (\Z_2)^n$.
    Also besitzt \(T^n\) genau $2^n$ verschiedene Spin-Strukturen.

    \item \textbf{\(2^{2g}\) Spin-Strukturen auf einer orientierten Fläche vom Geschlecht \(g\).}  
    Ist \(\Sigma_g\) eine geschlossene orientierte Fläche vom Geschlecht \(g\), so ist \(\Sigma_g\) spin und $H^1(\Sigma_g;\Z_2)\cong (\Z_2)^{2g}$. Daher besitzt \(\Sigma_g\) genau $2^{2g}$ verschiedene Spin-Strukturen. Insbesondere hat der Torus \(T^2=\Sigma_1\) genau \(4\) Spin-Strukturen.

\end{enumerate}

\end{EXA}


\bigskip



\section{Natürlicher Zusammenhang}

Wir wollen jetzt sicherstellen, dass wir den \emph{Spinorzusammenhang} als natürlichen Zusammenhang auf dem Spinorbündel aus den Riemannschen Daten ableiten können. Das läuft über Standardverfahren der Zusammenhangs-Theorie auf Bündeln. Dabei nutzen wir für Prinzipalbündel Zusammenhangsformen als verwandte Strukturen zu Zusammenhängen: 


\begin{DEF}{}{Zusammenhangsform auf $G$-Prinzipalbündeln}
  Sei $\pi:P\to M$ ein $G$-Prinzipalbündel. Eine Zusammenhangsform auf einem $G$-Prinzipalbündel ist eine $\mathfrak{g}$-wertige 1-Form
  $$
  \omega \in \Omega^1(P, \mathfrak{g}),
  $$
  für die gilt:
  \begin{enumerate}
    \item Reproduktion der Fundamentalfelder, d.h. für Fundamentalfelder von $\xi\in \mathfrak{g}$
    $$\xi_p^{\#}:=\left.\frac{d}{d t}\right|_{t=0}(p \cdot \exp (t \xi)) \in T_p P$$
    gilt
    $$
    \omega\left(\xi^{\#}\right)=\xi \quad \forall \xi \in \mathfrak{g} .
    $$
    \item Äquivarianz unter Rechtswirkung
    $$
    \left(R_g\right)^* \omega=\operatorname{Ad}\left(g^{-1}\right) \omega \quad \forall g \in G .
    $$
    Punktweise bedeutet das:
    $$
    \omega_{p \cdot g}\left(\left(R_g\right)_* v\right)=\operatorname{Ad}\left(g^{-1}\right)\left(\omega_p(v)\right) \quad \forall v \in T_p P .
    $$
  \end{enumerate}
\end{DEF}


\begin{REM}{}{Zusammenhangsform auf $\SO(n)$-Rahmenbündel}
  Sei $P^{\SO}(M)$ das $\SO(n)$-Rahmenbündel zu einer orientierten Riemannschen Mannigfaltigkeit $(M,g,\nabla^{LC})$ mit Levi-Civita Zusammenhang. Dann wird durch den Levi-Civita Zusammenhang des Tangentialbündels eine eindeutige Zusammenhangsform 
  $$\widetilde\omega^{LC}\in \Omega^1(P^\SO,\mathfrak{so}(n))$$ 
  auf $P^\SO(M)$ induziert.
\end{REM}



\begin{PROP}{}{Zusammenhangsform auf Spin-Prinzipalbündel}
  Sei $(P^\Spin(M),\Theta)$ die Spin-Struktur eine Spin-Mannigfaltigkeit. Dann wird durch die Levi-Civita Zusammenhangform $\widetilde\omega^{LC}\in\Omega^1(P^\SO(M),\mathfrak{so}(n))$ eine eindeutige Zusammenhangsform 
  $$\omega^{LC}\in \Omega^1(P^{\Spin}(M),\mathfrak{spin}(n))$$ 
  auf $P^\Spin(M)$ induziert, indem man die Zusammenhangsform $\widetilde\omega^{LC}$ längs der zweiblättrigen Überlagerung $\Theta: P^\Spin(M) \to P^\SO(M)$ pullbackt 
  $${\Theta}^*\widetilde\omega^{LC}\in \Omega^1(P^\Spin(M),\mathfrak{so}(n))$$
  und den Wertebereiches durch Anwendung des Liealgebren Isomorphismus $\lambda_*^{-1}:\mathfrak{so}(n)\to\mathfrak{spin}(n)$ anpasst
  $$\omega^{LC}=\lambda_*^{-1}\circ \left(\Theta^*\widetilde \omega^{LC}\right)\in \Omega^1(P^\Spin(M),\mathfrak{spin}(n)).$$

\end{PROP}

\begin{proof}
Wir weisen die Reproduzierbarkeit der Fundamentalfelder und Rechtsäquivarianz mittels Techniken aus \parencite[]{Kraus2025}:
  \begin{enumerate}
    \item \emph{Reproduktion der Fundamentalvektorfelder}: Für $\xi \in \mathfrak{spin}(n)$, $p\in P^\Spin(M)$ gilt:
    $$
    \xi_p^{\#}=\left.\frac{d}{d t}\right|_0 p \cdot \exp (t \xi) .
    $$
    Wegen der $\lambda$-Äquivarianz von $\Theta$ folgt
    $$
    d \Theta_p\left(\xi_p^{\#}\right)=\left.\frac{d}{d t}\right|_0 \Theta(p \cdot \exp (t \xi))=\left.\frac{d}{d t}\right|_0 \Theta(p) \cdot \lambda(\exp (t \xi))=\left.\frac{d}{d t}\right|_0 \Theta(p) \cdot \exp (t \lambda_* \xi)=\left(\lambda_* \xi\right)_{\Theta(p)}^{\#} .
    $$
    Daher
    $$
    \omega_p^{L C}\left(\xi_p^{\#}\right)=\lambda_*^{-1}\left(\widetilde{\omega}_{\Theta(p)}^{L C}\left(d \Theta_p\left(\xi_p^{\#}\right)\right)\right)=\lambda_*^{-1}\left(\widetilde{\omega}^{L C}\left(\left(\lambda_* \xi\right)^{\#}\right)\right)=\lambda_*^{-1}\left(\lambda_* \xi\right)=\xi .
    $$

    \item \emph{Rechtsäquivarianz}: Für $a \in \operatorname{Spin}(n)$ gilt
    $$\Theta \circ R_a=R_{\lambda(a)} \circ \Theta.$$
    Also
    $$
    \left(R_a\right)^* \omega^{L C}
    =\lambda_*^{-1} \circ \left(\left(R_a\right)^*\left(\Theta^* \tilde{\omega}^{L C}\right)\right)
    =\lambda_*^{-1} \circ \left(\left(\Theta \circ R_a\right)^* \tilde{\omega}^{L C}\right)
    =\lambda_*^{-1} \circ\left(\left(R_{\lambda(a)} \circ \Theta\right)^* \tilde{\omega}^{L C}\right).
    $$
    Da $\widetilde{\omega}^{L C}$ eine Zusammenhangsform auf $P^{S O}(M)$ ist, gilt
    $$
    \left(R_{\lambda(a)}\right)^* \widetilde{\omega}^{L C}= \operatorname{Ad}\left(\lambda(a)^{-1}\right) \widetilde{\omega}^{L C} .
    $$
    Somit
    $$
    \left(R_a\right)^* \omega^{L C}
    =\lambda_*^{-1} \circ \left(\Theta^*(R_{\lambda(a)}^* \tilde{\omega}^{L C})\right)
    =\lambda_*^{-1} \circ \left(\Theta^* \left(\operatorname{Ad}\left(\lambda(a)^{-1}\right) \widetilde{\omega}^{L C}\right)\right)
    =\lambda_*^{-1} \circ \left( \operatorname{Ad} (\lambda(a)^{-1}) \Theta^* \widetilde{\omega}^{L C}\right).
    $$
    Wir wollen jetzt
    $$
    \lambda_* \circ  \operatorname{Ad}(a)= \operatorname{Ad}(\lambda(a)) \circ \lambda_* .
    $$
    nutzen. Man kann das bspw. kategorientheoretisch einsehen (für einen kategorientheoretischen Zugang, siehe \parencite[p.~290]{HilgertNeeb2012}).
    Für die Konjugation $c_a$ mit $a\in\Spin(n)$ und der zweifachen Überlagerung $\lambda$ von $\SO(n)$ durch $\Spin(n)$ gilt
    $$\lambda \circ c_a=c_{\lambda(a)} \circ \lambda.$$
    Dann folgt mit der Funktorialität des \emph{Lie-Funktor} $L$ :
    $$
    L(\lambda) \circ L\left(c_a\right)=L\left(\lambda \circ c_a\right)=L\left(c_{\lambda(a)} \circ \lambda\right)=L\left(c_{\lambda(a)}\right) \circ L(\lambda) .
    $$
    Da
    $$
    L(\lambda)=\lambda_*, \quad L\left(c_a\right)= \operatorname{Ad}(a), \quad L\left(c_{\lambda(a)}\right)= \operatorname{Ad}(\lambda(a)),
    $$
    erhält man
    $$
    \lambda_* \circ  \operatorname{Ad}(a)= \operatorname{Ad}(\lambda(a)) \circ \lambda_*.
    $$
    Anwenden für $a\to a^{-1}$ gibt also
    $$
    \left(R_a\right)^* \omega^{L C}= \operatorname{Ad}\left(a^{-1}\right) \omega^{L C} .
    $$
  \end{enumerate} 
\end{proof}

\bigskip

Damit können wir den Spinorzusammenhang konstruieren. 

\begin{PROP}{}{Kovariante Ableitung auf dem Spinorbündel}
Sei $\Sigma M= P^\Spin(M)\times_{\sigma_n}\Sigma_n$ das Spinorbündel einer Spin-Mannigfaltigkeit \((M,g,P^\Spin(M))\), und sei $\omega^{LC}\in \Omega^1(P^\Spin(M),\mathfrak{spin}(n))$ die von der Levi-Civita-Verbindung induzierte Zusammenhangsform auf dem Spin-Prinzipalbündel. Dann induziert \(\omega^{LC}\) eine kovariante Ableitung
\[
\nabla^\Sigma:\Gamma(TM)\times \Gamma(\Sigma M)\to \Gamma(\Sigma M),
\]
die lokal wie folgt beschrieben wird: Ist \(U\subset M\) offen, \(H\in \Gamma(P^\Spin(M)|_U)\) ein lokaler Spin-Rahmen und
\[
s=[H,\varphi]\in \Gamma(\Sigma M|_U)
\qquad\text{mit}\qquad
\varphi\in C^\infty(U,\Sigma_n),
\]
so gilt für jedes \(X\in \Gamma(TM|_U)\)
\[
\nabla_X^\Sigma s
=
\left[H,\partial_X\varphi+(\sigma_n)_*\bigl(H^*\omega^{LC}(X)\bigr)\varphi\right].
\]
Hierbei ist $H^*\omega^{LC}\in \Omega^1(U,\mathfrak{spin}(n))$ die lokale Zusammenhangs-\(1\)-Form und $(\sigma_n)_*:\mathfrak{spin}(n)\to \End(\Sigma_n)$ die abgeleitete Darstellung.
\end{PROP}

\begin{proof}
Wir zeigen zunächst, dass die angegebene Formel wohldefiniert ist, also nicht von der Wahl des lokalen Spin-Rahmens \(H\) abhängt.

\medskip

\begin{enumerate}
    \item \emph{Wohldefiniertheit.}
    Sei \(a:U\to \Spin(n)\) glatt und
    \[
    H':=H\cdot a.
    \]
    Dann haben wir äquivalente lokale Spinorschnitt \(s\)
    \[
    s= [h,\varphi]= [H',\varphi']
    \qquad\text{mit}\qquad
    \varphi':=\sigma_n(a^{-1})\varphi.
    \]
    Setze
    \[
    \theta:=H^*\omega^{LC}\in \Omega^1(U,\mathfrak{spin}(n)),
    \qquad
    \theta':=(H')^*\omega^{LC}\in \Omega^1(U,\mathfrak{spin}(n)).
    \]
    Für den Wechsel des lokalen Schnittes gilt die Transformationsformel \parencite[38]{Kraus2025}
    \[
    \theta'=\Ad(a^{-1})\,\theta+a^{-1}da.
    \]
    Weiter gilt für die abgeleitete Darstellung wie oben durch $\sigma_n\circ c_a=c_{\sigma_n(a)}\circ\sigma_n$
    \[
    (\sigma_n)_*(\Ad(a^{-1})\xi)
    =
    \sigma_n(a^{-1})\,(\sigma_n)_*(\xi)\,\sigma_n(a)
    \qquad
    \xi\in \mathfrak{spin}(n).
    \]
    Außerdem erhält man für die Ableitung von \(\varphi'=\sigma_n(a^{-1})\varphi\):
    \[
    \partial_X\varphi'
    =
    \partial_X\bigl(\sigma_n(a^{-1})\varphi\bigr)
    =
    \bigl(\partial_X(\sigma_n(a^{-1}))\bigr)\varphi+\sigma_n(a^{-1})\partial_X\varphi.
    \]

    Wir müssen $\partial_X(\sigma_n(a^{-1}))$ bestimmen. 
    Setze dazu $x \in U$ fest und wähle eine Kurve $c(t)$ in $U$ mit
    $$c(0)=x, \quad \dot{c}(0)=X_x.$$
    Dann haben wir eine Kurve in $\Spin(n)$ durch $g(t):=a(c(t)) \subset  \Spin(n)$ mit
    $$g(0)=a(x),\quad g'(0)=da_x(X_x)\in T_{a(x)}\Spin(n)$$
    Durch die linke Maurer--Cartan-Form erhalten wir ein Element von
    \(\mathfrak{spin}(n)\):
    \[
    (a^{-1}da)(X_x)
    =
    \theta^L_{a(x)}\bigl(da_x(X_x)\bigr)
    =
    (dL_{a(x)^{-1}})_{a(x)}\bigl(da_x(X_x)\bigr)
    \in \mathfrak{spin}(n).
    \]
    Schreibt man \(g(t):=a(c(t))\), so ist dies in der Clifford-Algebra-Schreibweise gerade
    \[
    (a^{-1}da)(X_x)=g(0)^{-1}g'(0).
    \]
    Nun schreiben wir
    $$h(t):=g(0)^{-1} g(t), \quad \text{also} \quad g(t)=g(0) h(t).$$
    Dann gilt
    $$h(0)=e, \quad h^{\prime}(0)=g(0)^{-1} g^{\prime}(0)=\left(a^{-1} d a\right)\left(X_x\right).$$
    Wegen der Homomorphieeigenschaft der Darstellung folgt
    $$\sigma_n(g(t))=\sigma_n(g(0)) \sigma_n(h(t)).$$
    Differenzieren bei $t=0$ ergibt
    $$\left.\frac{d}{d t}\right|_0 \sigma_n(g(t))=\left.\sigma_n(g(0)) \frac{d}{d t}\right|_0 \sigma_n(h(t)).$$
    Da $h(0)=e$, ist
    $$\left.\frac{d}{d t}\right|_0 \sigma_n(h(t))=\left(\sigma_n\right)_*\left(h^{\prime}(0)\right)=\left(\sigma_n\right)_*\left(\left(a^{-1} d a\right)\left(X_x\right)\right)$$
    Also
    $$\partial_X\left(\sigma_n(a)\right)=\sigma_n(a)\left(\sigma_n\right)_*\left(\left(a^{-1} d a\right)(X)\right).$$
    Durch Ableiten mit Produktregel erhält man aus 
    $$\sigma_n(a) \sigma_n\left(a^{-1}\right)=\mathrm{id}_{\Sigma_n}$$
    abschließend 
    $$\partial_X\left(\sigma_n\left(a^{-1}\right)\right)=-\left(\sigma_n\right)_*\left(\left(a^{-1} d a\right)(X)\right) \sigma_n\left(a^{-1}\right).$$

    Damit können wir jetzt die Ableitung bestimmen:
    \[
    \partial_X\varphi'
    =
    -(\sigma_n)_*\bigl((a^{-1}da)(X)\bigr)\sigma_n(a^{-1})\varphi
    +\sigma_n(a^{-1})\partial_X\varphi.
    \]
    Es folgt
    \begin{align*}
    \partial_X\varphi'+(\sigma_n)_*(\theta'(X))\varphi'
    &=
    -(\sigma_n)_*\bigl((a^{-1}da)(X)\bigr)\sigma_n(a^{-1})\varphi
    +\sigma_n(a^{-1})\partial_X\varphi \\
    &\quad
    +(\sigma_n)_*\bigl(\Ad(a^{-1})\theta(X)\bigr)\sigma_n(a^{-1})\varphi
    +(\sigma_n)_*\bigl((a^{-1}da)(X)\bigr)\sigma_n(a^{-1})\varphi \\
    &=
    \sigma_n(a^{-1})\partial_X\varphi
    +\sigma_n(a^{-1})(\sigma_n)_*(\theta(X))\varphi \\
    &=
    \sigma_n(a^{-1})
    \bigl(\partial_X\varphi+(\sigma_n)_*(\theta(X))\varphi\bigr).
    \end{align*}
    Also ist
    \begin{align*}
    \left[H',\partial_X\varphi'+(\sigma_n)_*(\theta'(X))\varphi'\right]
    &=
    \left[H\cdot a,\sigma_n(a^{-1})
    \bigl(\partial_X\varphi+(\sigma_n)_*(\theta(X))\varphi\bigr)\right] \\
    &=
    \left[H,\partial_X\varphi+(\sigma_n)_*(\theta(X))\varphi\right].
    \end{align*}
    Damit ist die Definition unabhängig von der Wahl des lokalen Spin-Rahmens und somit wohldefiniert.


    \item \emph{Linearität in der Vektorfeldvariable.}
    Seien \(X,Y\in \Gamma(TU)\) und \(f\in C^\infty(U)\). Dann gilt unmittelbar aus der Formel
    \begin{align*}
    \nabla^\Sigma_{X+Y}[H,\varphi]
    &=
    \left[H,\partial_{X+Y}\varphi+(\sigma_n)_*(\theta(X+Y))\varphi\right] \\
    &=
    \left[H,\partial_X\varphi+(\sigma_n)_*(\theta(X))\varphi\right]
    +
    \left[H,\partial_Y\varphi+(\sigma_n)_*(\theta(Y))\varphi\right] \\
    &=
    \nabla^\Sigma_X[H,\varphi]+\nabla^\Sigma_Y[H,\varphi],
    \end{align*}
    sowie
    \begin{align*}
    \nabla^\Sigma_{fX}[H,\varphi]
    &=
    \left[H,\partial_{fX}\varphi+(\sigma_n)_*(\theta(fX))\varphi\right] \\
    &=
    \left[H,f\,\partial_X\varphi+f\,(\sigma_n)_*(\theta(X))\varphi\right] \\
    &=
    f\,\nabla^\Sigma_X[H,\varphi].
    \end{align*}
    Also ist \(\nabla^\Sigma\) \(C^\infty(M)\)-linear in der Vektorfeldvariable.

    \medskip

    \item \emph{Additivität und Leibniz-Regel in der Schnittvariable.}
    Für lokale Spinorfelder \(s_1=[H,\varphi_1]\), \(s_2=[H,\varphi_2]\) gilt
    \begin{align*}
    \nabla^\Sigma_X(s_1+s_2)
    &=
    \nabla^\Sigma_X[H,\varphi_1+\varphi_2] \\
    &=
    \left[H,\partial_X(\varphi_1+\varphi_2)+(\sigma_n)_*(\theta(X))(\varphi_1+\varphi_2)\right] \\
    &=
    \nabla^\Sigma_X s_1+\nabla^\Sigma_X s_2.
    \end{align*}
    Ferner gilt für \(f\in C^\infty(U)\):
    \begin{align*}
    \nabla^\Sigma_X(f[H,\varphi])
    &=
    \nabla^\Sigma_X[H,f\varphi] \\
    &=
    \left[H,\partial_X(f\varphi)+(\sigma_n)_*(\theta(X))(f\varphi)\right] \\
    &=
    \left[H,X(f)\varphi+f\,\partial_X\varphi+f\,(\sigma_n)_*(\theta(X))\varphi\right] \\
    &=
    X(f)[H,\varphi]+f\left[H,\partial_X\varphi+(\sigma_n)_*(\theta(X))\varphi\right] \\
    &=
    X(f)\,s+f\,\nabla^\Sigma_X s.
    \end{align*}
    Damit ist die Leibniz-Regel erfüllt.
\end{enumerate}


Da die lokalen Formeln auf Überlappungen wegen Schritt 1 übereinstimmen, verkleben sie zu einer global wohldefinierten kovarianten Ableitung auf \(\Sigma M\).
\end{proof}

\begin{REM}{}{Eigenschaften des Spinorzusammenhangs}
  \begin{enumerate}
    \item \emph{Christoffel-Symbole} \parencite[101]{BaerSpin}: Die kovariante Ableitung von lokalen Spinorschnitten $[H, \varphi ] \in \Gamma(\Sigma M|_U)$ kann durch Christoffel-Symbole der Levi-Civita Zusammenhangsform $\widetilde\omega^{LC}$ des $\SO(n)$-Rahmenbündels dargestellt werden: Bzgl. des induzierten ON-Rahmen $b_i=h(e_i)=[h,e_i]\in\mathfrak{X}(U)$ zu $h=\Theta\circ H:U\to P ^\SO(M)$ gilt
    $$
    \begin{aligned}
    \nabla_{b_i}^{\Sigma} [ H, \varphi ] = [ H, \partial_{b_i} \varphi+\frac{1}{4} \sum_{j, k=1}^n \Gamma_{i j}^k e_j \cdot e_k \cdot \varphi ]
    \end{aligned}
    $$
    wobei die Christoffel-Symbole bestimmt werden durch 
    $$\Gamma_{i j}^k=\widetilde\omega_{k j}^{\mathrm{LC}}\left(d h\left(b_i\right)\right).$$
    \item \emph{Leibniz-Regel bzgl. Clifford-Multiplikation} \parencite[103]{BaerSpin}:  Der Spinorzusammenhang erfüllt eine Leibniz-Regel für Clifford-Multiplikation, d.h. für $s\in\Gamma(\Sigma M)$ und $X, Y \in \mathfrak{X}(M)$ gilt
    $$
    \nabla_X^{\Sigma}(Y \cdot s)=\left(\nabla_X^{\mathrm{LC}} Y\right) \cdot s+Y \cdot \nabla_X^{\Sigma} s.
    $$
  \end{enumerate}
\end{REM}


\begin{DEF}{}{Krümmung des Spinorzusammenhangs}
  Die Krümmung des Spinorzusammenhangs (Krümmungsendomorphismus) ist für $X,Y\in\mathfrak{X}(M)$ gegeben durch 
    $$R^{\Sigma}(X, Y): \Sigma M \rightarrow \Sigma M;\quad R^{\Sigma}(X, Y) s:=\nabla_X^{\Sigma} \nabla_Y^{\Sigma} s-\nabla_Y^{\Sigma} \nabla_X^{\Sigma} s-\nabla_{[X, Y]}^{\Sigma} s$$
\end{DEF}


\begin{REM}{\parencite[104,108]{BaerSpin}}{Identitäten der Krümmung des Spinorzusammenhangs}
  Sei $(M,g,\nabla^{LC})$ eine Riemannsche Spin-Mannigfaltigkeit und $(\Sigma M,\nabla^\Sigma)$ das zugehörige Spinorbündel. Sei ferner $X,Y\in T_x M$, $b_1,...,b_n$ eine ONB von $T_xM$ und $\varphi\in \Sigma_n$ ein Spinor. Dann gelten punktweisen die Relationen:
  $$
  \begin{aligned}
    &R^{\Sigma}(X, Y) \varphi =-\frac{1}{4} \sum_{i=1}^n\left(R(X, Y) b_i\right) \cdot b_i \cdot \varphi\\
    &\sum_{j=1}^n b_j \cdot R^{\Sigma}\left(b_j, X\right) \varphi=\frac{1}{2} \operatorname{Ric}(X) \cdot \varphi
  \end{aligned}
  $$
  wobei $R$ die Riemannscher Krümmung und $\operatorname{Ric}$ der Ricci-Endomorphismus des Levi-Civita Zusammenhangs ist. 
\end{REM}



\section{Natürliche Hermitesche Bündelmetrik auf Spinorbündeln}



\begin{DEF}{}{Hermitesches Skalarprodukt auf dem Spinorraum}
Wir definieren auf $\Sigma_n$ das \emph{Hermitesche Skalarprodukt}
\[
\langle\cdot,\cdot\rangle_\Sigma:\Sigma_n\times \Sigma_n\to \C
\]
durch die Vorschrift 
$$\langle\varphi, \psi\rangle_{\Sigma}:=\sum_J a_J \overline{b_J} \quad \text{für} \quad \varphi=\sum_J a_J z(J)\cdot\omega,\ \psi=\sum_J b_J z(J)\cdot\omega$$
als die eindeutig bestimmte Abbildung, die im ersten Argument komplex-linear, im zweiten Argument komplex-antilinear ist \emph{und} für die die angegebene Basis $\{z(J_k)\omega\}_J$ des Spinorraums eine \emph{Orthonormalbasis} bilden, d.\,h.
\[
\bigl\langle z(J_k)\omega,\, z(I_\ell)\omega\bigr\rangle
=
\delta_{k\ell}\,\delta_{(j_1,\dots,j_k),(i_1,\dots,i_\ell)}.
\]
\end{DEF}

\begin{REM}{}{Eigenschaften des hermiteschen Skalarprodukt auf $\Sigma_n$}
  \begin{enumerate}
    \item \emph{Hermitesche Symmetrie}:
    $$\langle\varphi, \psi\rangle_{\Sigma}=\overline{\langle\psi, \varphi\rangle_{\Sigma}}$$
    \item \emph{Positivität}:
    $$
    \langle\phi, \phi\rangle_\Sigma=\sum_I\left|a_I\right|^2 \geq 0,
    $$
    und Gleichheit gilt genau dann, wenn $\phi=0$.
    \item \emph{Clifford-Multiplikation schiefadjungiert} \parencite[89]{BaerSpin}: Die Clifford-Multiplikation von Spinoren mit Vektoren $X\in\R^n$ ist schiefadjungiert bzgl $\langle\cdot,\cdot\rangle_\Sigma$:
    $$\langle X \cdot \phi, \psi\rangle_\Sigma= -\langle\phi, X \cdot  \psi\rangle_\Sigma$$
    \item \emph{Orthogonalität der Chiralitätszerlegung}: Da $\Sigma_n^{+}$und $\Sigma_n^{-}$von disjunkten Teilmengen dieser Orthonormalbasis erzeugt werden, ist
    $$
    \Sigma_n=\Sigma_n^{+} \oplus \Sigma_n^{-}
    $$
    eine orthogonale Zerlegung.
  \end{enumerate}
  
\end{REM}




\begin{DEF}{}{Kanonische Hermitesche Bündelmetrik auf dem Spinorbündel}
Sei \(M\) eine Riemannsche Spin-Mannigfaltigkeit mit Spin-Struktur $P^{\Spin}(M)\to M$, und sei $\Sigma M=P^{\Spin}(M)\times_{\sigma_n}\Sigma_n$
das zugehörige Spinorbündel. Auf \(\Sigma M\) definieren wir faserweise eine
Hermitesche Bündelmetrik $\langle\cdot,\cdot\rangle_{\Sigma M}$ durch
\[
\bigl\langle [H,\varphi],[H,\psi]\bigr\rangle_{\Sigma M}
:=
\langle \varphi,\psi\rangle_{\Sigma},
\]
wobei faserweise durch $\Sigma M|_x\cong\Sigma_n$ das hermiteschen Skalarprodukts $(\Sigma_n,\langle\cdot,\cdot\rangle_\Sigma)$ verwendet wird. Diese heißt die
\emph{kanonische Hermitesche Bündelmetrik} auf \(\Sigma M\).
\end{DEF}

\begin{PROP}{}{Wohldefiniertheit der kanonischen Hermiteschen Bündelmetrik}
Die obige Vorschrift definiert eine wohldefinierte Hermitesche Bündelmetrik auf
dem Spinorbündel \(\Sigma M\).
\end{PROP}

\begin{proof}
Da \(\Sigma M=P^{\Spin}(M)\times_{\sigma_n}\Sigma_n\) ein assoziiertes
Vektorbündel ist, gilt für alle \(a\in \Spin(n)\)
\[
[H,\varphi]=[H\cdot a,\sigma_n(a^{-1})\varphi].
\]
Um die Wohldefiniertheit zu zeigen, genügt es daher nachzuprüfen, dass die
definierte Form unabhängig von der Wahl des Repräsentanten ist.

Seien also \(a\in \Spin(n)\), \(H\in P^{\Spin}(M)\) und \(\varphi,\psi\in \Sigma_n\).
Dann gilt
\[
\begin{aligned}
\bigl\langle [H\cdot a,\sigma_n(a^{-1})\varphi],
[H\cdot a,\sigma_n(a^{-1})\psi]\bigr\rangle_{\Sigma M}
&=
\bigl\langle \sigma_n(a^{-1})\varphi,\sigma_n(a^{-1})\psi\bigr\rangle_{\Sigma_n}.
\end{aligned}
\]
Da die Spinor-Darstellung \(\sigma_n\) unitär bezüglich
\(\langle\cdot,\cdot\rangle_{\Sigma_n}\) ist, folgt
\[
\bigl\langle \sigma_n(a^{-1})\varphi,\sigma_n(a^{-1})\psi\bigr\rangle_{\Sigma_n}
=
\langle \varphi,\psi\rangle_{\Sigma_n}.
\]
Also ist
\[
\bigl\langle [H\cdot a,\sigma_n(a^{-1})\varphi],
[H\cdot a,\sigma_n(a^{-1})\psi]\bigr\rangle_{\Sigma M}
=
\bigl\langle [H,\varphi],[H,\psi]\bigr\rangle_{\Sigma M},
\]
und die Vorschrift ist wohldefiniert.

Sesquilinearität, Hermitesche Symmetrie und positive Definitheit folgen jeweils
unmittelbar faserweise aus den entsprechenden Eigenschaften des Hermiteschen
Skalarprodukts auf \(\Sigma_n\). Somit ist \(\langle\cdot,\cdot\rangle_{\Sigma M}\)
eine Hermitesche Bündelmetrik auf \(\Sigma M\).
\end{proof}




\begin{REM}{}{Eigenschaften des Spinorzusammenhangs und Clifford-Multiplikation unter Skalarprodukt}
    \begin{enumerate}
        \item \emph{Spinorzusammenhang ist Metrik-Zusammenhang} \parencite[102]{BaerSpin}: Der Spinorzusammenhang $\nabla^\Sigma$ ist verträgt sich mit der Metrik, d.h. für Spinorfelder $s,s'\in\Gamma(\Sigma M)$ und Vektorfelder $X\in\mathfrak{X}(M)$ gilt:
        $$
        \partial_X\langle s, s'\rangle_{\Sigma M}=\left\langle\nabla_X^{\Sigma} s, s'\right\rangle_{\Sigma M}+\left\langle s, \nabla_X^{\Sigma} s'\right\rangle _{\Sigma M}.
        $$
        \item \emph{Schiefadjungiertheit der Clifford-Multiplikation} \parencite[103]{BaerSpin}: Bezüglich der kanonischen Hermiteschen Bündelmetrik auf \(\Sigma M\) gilt
        punktweise
        \[ \langle X\cdot \phi,\psi\rangle_{\Sigma M} = -\langle \phi, X\cdot \psi\rangle_{\Sigma M}\]
        für alle \(X\in T_xM\) und \(\phi,\psi\in \Sigma_xM\). Äquivalent dazu ist
        \[
        c(X)^*=-c(X)\in \End(\Sigma M_x).
        \]
    \end{enumerate}
\end{REM}





\section{Dirac Operator}


\begin{DEF}{}{Dirac Operator}
Sei $(M,g)$ eine Riemannsche Spin-Mannigfaltigkeit, $\Sigma M$ das zugehörige
Spinorbündel und $\nabla^\Sigma$ der vom Levi-Civita-Zusammenhang induzierte
Spinorzusammenhang. Sei ferner die Clifford-Multiplikation
\[
c:T^*M\otimes \Sigma M\to \Sigma M,\qquad c(\alpha\otimes \varphi):=\alpha^\sharp\cdot\varphi,
\]
als eine Abbildung auf Schnitten
\[
c:\Gamma(T^*M\otimes \Sigma M)\to \Gamma(\Sigma M).
\]
aufgefasst. Der \emph{Dirac-Operator} ist der Differentialoperator erster Ordnung
\[
D:=c\circ \nabla^\Sigma:\Gamma(\Sigma M)\to \Gamma(\Sigma M).
\]
Bezüglich eines lokalen orthonormalen Rahmens $b_1,\dots,b_n$ über $M$ mit dualem
Korahmen $b_1^*,\dots,b_n^*$ gilt lokal
\[
D\varphi
=
\sum_{i=1}^n c\bigl(b_i^*\otimes \nabla^\Sigma_{b_i}\varphi\bigr)
=
\sum_{i=1}^n (b_i^*)^\sharp\cdot \nabla^\Sigma_{b_i}\varphi
=
\sum_{i=1}^n b_i\cdot \nabla^\Sigma_{b_i}\varphi.
\]
\end{DEF}


\begin{PROP}{}{Formale Selbstadjungiertheit des Dirac-Operators} 
  Der Dirac-Operator ist \emph{formal selbstadjungiert}, d.h. für alle $\varphi,\psi \in \Gamma(\Sigma M)$ mit kompaktem Träger gilt
  $$
  (D\varphi,\psi)_{\Sigma M}=\int_M\langle D \varphi, \psi\rangle_{\Sigma M} dvol = \int_M\left\langle \varphi, D \psi\right\rangle_{\Sigma M} dvol= (\varphi,D\psi)_{\Sigma M}
  $$
\end{PROP}

\begin{proof}
Für jedes $x\in M$ definiert
\[
\alpha_x:T_xM\to\mathbb C,\qquad
\alpha_x(Y):=\langle Y\cdot\phi(x),\psi(x)\rangle_{\Sigma_xM}.
\]
Da die Clifford-Multiplikation linear in $Y$ ist, ist $\alpha_x$ ein
komplex-wertiges lineares Funktional auf $T_xM$, also
\[
\alpha_x\in \Hom(T_xM,\mathbb C)\cong T_x^*M\otimes\mathbb C.
\]
Wegen der Glattheit der Clifford-Multiplikation und des hermiteschen
Faserprodukts der Spinoren erhält man so einen glatten Schnitt
\[
\alpha\in \Gamma(T^*M\otimes\mathbb C),\qquad
\alpha(Y)=\langle Y\cdot\phi,\psi\rangle_{\Sigma M}.
\]
Über den durch die Riemannsche Metrik induzierten musikalischen Isomorphismus
\[
\sharp:T^*M\otimes\mathbb C\to TM\otimes\mathbb C
\]
gibt es genau ein Vektorfeld
\[
X:=\alpha^\sharp\in \Gamma(TM\otimes\mathbb C)
\]
mit
\[
\langle X,Y\rangle=\alpha(Y)=\langle Y\cdot\phi,\psi\rangle
\qquad\text{für alle }Y\in \mathfrak X(M),
\]
wobei $\langle\cdot,\cdot\rangle$ links die komplex-bilineare Fortsetzung der
Riemannschen Metrik auf $TM\otimes\mathbb C$ bezeichnet.
Da $\phi,\psi\in \Gamma(\Sigma M)$ kompakt getragen sind, ist auch
$X\in \Gamma(TM\otimes\mathbb C)$ kompakt getragen.

Sei $b_1, \ldots, b_n$ ein lokaler orthogonaler Rahmen. Dann gilt:
$$
  \begin{aligned}
  \operatorname{div} X & =\sum_{i=1}^n\left\langle\nabla_{b_i} X, b_i\right\rangle \\
  & =\sum_{i=1}^n \partial_{b_i}\left\langle X, b_i\right\rangle-\left\langle X, \nabla_{b_i} b_i\right\rangle \\
  & =\sum_{i=1}^n \partial_{b_i}\left\langle b_i \cdot \phi, \psi\right\rangle_{\Sigma M}-\left\langle\nabla_{b_i} b_i \cdot \phi, \psi\right\rangle_{\Sigma M} \\
  & = \sum_{i=1}^n\left\langle\nabla_{b_i}^{\Sigma}\left(b_i \cdot \phi\right), \psi\right\rangle_{\Sigma M}+\left\langle b_i \cdot \phi, \nabla_{b_i}^{\Sigma} \psi\right\rangle_{\Sigma M}-\left\langle\nabla_{b_i} b_i \cdot \phi, \psi\right\rangle_{\Sigma M}\\
  & = \sum_{i=1}^n\left\langle\nabla_{b_i} b_i \cdot \phi+b_i \cdot \nabla_{b_i}^{\Sigma} \phi, \psi\right\rangle_{\Sigma M}-\left\langle\phi, b_i \cdot \nabla_{b_i}^{\Sigma} \psi\right\rangle_{\Sigma M}-\left\langle\nabla_{b_i} b_i \cdot \phi, \psi\right\rangle_{\Sigma M} \\ 
  & = \langle D \phi, \psi\rangle_{\Sigma M}-\langle\phi, D \psi\rangle_{\Sigma M}
  \end{aligned}
  $$
  Da alle Ausdrücke komplexwertig sind, wenden wir den Divergenzsatz komponentenweise (Real-/Imaginärteil) an:
  $$
  (D \phi, \psi)_{\Sigma M}-(\phi, D \psi)_{\Sigma M}=\int_M \operatorname{div}(X) d v o l=0
  $$
\end{proof}

\begin{REM}{}{Eigenschaften des Dirac-Operators}
  \begin{enumerate}
    \item \emph{Prinzipalsymbol}: Das Prinzipalsymbol des Dirac-Operators ist die \emph{Clifford-Multiplikation}:
    $$\sigma_1(D, \xi) \varphi=\sigma_1\left(D_{A, \nabla^{\Sigma}}, \xi\right) \varphi = A(\xi \otimes \varphi)=\xi^{\sharp} \cdot \varphi=c(\xi\otimes \varphi)$$
    \item \emph{Dirac-Typ}: Der Dirac-Operator ist vom \emph{Dirac-Typ}, was für formal selbstadjungierten Differentialoperatoren mit übereinstimmenden Definitions- und Wertebereich äquivalent ist zu $D^2$ ist vom \emph{Laplace-Typ}:
    $$\sigma_2(D^2, \xi)=-|\xi|^2 \cdot \operatorname{id}_{\Sigma M_x} \quad \text{ für alle } x \in M \text{ und } \xi \in T_x^* M$$
    Dies folgt direkt aus Standardeigenschaften des Prinzipalsymbols und der Clifford-Relation:
    $$\sigma_2(D^2,\xi)\varphi=\sigma_1(D,\xi)\circ\sigma_1(D,\xi)\varphi=\xi^\sharp\cdot(\xi^\sharp\cdot \varphi)=(\xi^\sharp\cdot\xi^\sharp)\cdot \varphi=-|\xi|_g^2\varphi$$


    
    
  \end{enumerate}
\end{REM}





\section{Lichnerowicz-Formel und Vanishing Theorem}






\begin{THEO}{}{Lichnerowicz (1963)}
Sei $(M,g)$ eine Riemannsche Spin-Mannigfaltigkeit, $\Sigma M$ das zugehörige
Spinorbündel und $D$ der Dirac-Operator. Dann gilt für die \emph{Dirac-Laplace-Operator} $D^2$ als dem quadrierten Dirac-Operator die Lichnerowicz-Formel
$$
D^2=\left(\nabla^{\Sigma}\right)^* \nabla^{\Sigma}+\frac{\operatorname{scal}}{4} \cdot \operatorname{id}_{\Sigma M}
$$  
\end{THEO}

\begin{proof}
  Wir geben eine Beweis-Skizze: Sei $x \in M$, und $b_1, \ldots, b_n$ ein lokaler ON-Rahmen in $TM$ um $x$, sodass 
  $$\nabla_{b_i}^{\mathrm{LC}} b_j(x)=0\qquad \text{ und }\qquad \operatorname{Ric}\left(b_i\right)(x)=\lambda_i b_i(x)$$
  für alle $i, j$ gilt. Für lokale Spinorfelder $\varphi \in \Gamma(\Sigma M)$ um $x$ gilt nach Anwendung von Leibniz und der parallelen transport des Rahmens bei $x$:
  $$\begin{aligned}
    \left(D^2 \phi\right)_{(x)} 
    & \stackrel{\text{Def}}{=}\left(\sum_{i, j=1}^n \left(b_j \cdot \nabla_{b_j}^{\Sigma}\right)\left( b_i \cdot \nabla_{b_i}^{\Sigma} \phi \right)\right)_{(x)} \\ 
    & \stackrel{\text{Leib}}{=}\left(\sum_{i, j=1}^n b_j \cdot(\nabla_{b_j}^{\mathrm{LC}} b_i \cdot \nabla_{b_i}^{\Sigma} \phi+b_i \cdot \nabla_{b_j}^{\Sigma} \nabla_{b_i}^{\Sigma} \phi)\right)_{(x)} \\ 
    & \stackrel{\text{PT}}{=}\left(\sum_{i, j=1}^n b_j \cdot b_i \cdot \nabla_{b_j}^{\Sigma} \nabla_{b_i}^{\Sigma} \phi)\right)_{(x)} \\ 
  \end{aligned}$$
  Spliten der Summe
  $$\sum_{i,j=1}^n\to \sum_{i=1}^n + \sum_{i<j}^n + \sum_{i>j}^n,$$
  Anwendung der Clifford-Relation für $ b_i \cdot b_i=-1$ und $b_j\cdot b_i=-b_i\cdot b_j$ und Anpassung der Indizes gibt:
  $$\begin{aligned} 
    & = \left(-\sum_{i=1}^n \nabla_{b_i}^{\Sigma} \nabla_{b_i}^{\Sigma} \phi+\sum_{i<j} b_j \cdot b_i \cdot\left(\nabla_{b_j}^{\Sigma} \nabla_{b_i}^{\Sigma}-\nabla_{b_i}^{\Sigma} \nabla_{b_j}^{\Sigma}\right) \phi\right)_{(x)} \\ 
    & =\left(\left(\nabla^{\Sigma}\right)^* \nabla^{\Sigma} \phi+\sum_{i<j} b_j \cdot b_i \cdot R^{\Sigma}\left(b_j, b_i\right) \phi\right)_{(x)}
  \end{aligned}$$
  In der ersten Summe identifizieren wir den Zusammenhangs-Laplace minus einen Summe über $\nabla_{\nabla_{b_i}^{LC}b_i}$ (siehe lokale Darstellung von Zusammenhangslaplace im Appendix), wobei die Summe nach Annahme in $x$ verschwindet, und in der zweiten Summe den Krümmungsendomorphismus $R^\Sigma$ des Spinorzusammenhangs plus eine Summe über $\nabla^\Sigma_{[b_j,b_i]}$-Terme, welche ebenfalls nach Annahme aus Torsionsfreiheit in $x$ verschwinden:
  $$\begin{aligned} 
    & = \left(\left(\nabla^{\Sigma}\right)^* \nabla^{\Sigma}\varphi - \sum_{i=1}^n\left(\nabla_{\nabla_{b_i}^{\mathrm{LC}} b_i}\right)\varphi +\sum_{i<j} b_j \cdot b_i \cdot\left(R^{\Sigma}\left(b_j, b_i\right) + \nabla^\Sigma_{[b_j,b_i]}\right) \phi\right)_{(x)} \\ 
    & =\left(\left(\nabla^{\Sigma}\right)^* \nabla^{\Sigma} \phi+\sum_{i<j} b_j \cdot b_i \cdot R^{\Sigma}\left(b_j, b_i\right) \phi\right)_{(x)}
  \end{aligned}$$
  Wir haben also schon den Zusammenhangs-Laplace und konzentrieren uns auf den zweiten Term. Zuerst bringen wir die Summe in eine geeignete Form. Dabei nutzen wir die Clifford-Relation und Antisymmetrie der Krümmung in ersten beiden Argumenten
  $$\begin{aligned}
    \left(\sum_{i<j} b_j \cdot b_i \cdot R^{\Sigma}\left(b_j, b_i\right) \phi\right)_{(x)} 
    & = \frac{1}{2}\left(\sum_{i<j} b_j \cdot b_i \cdot R^{\Sigma}\left(b_j, b_i\right) \phi + \sum_{i<j} b_i \cdot b_j \cdot R^{\Sigma}\left(b_i, b_j\right) \phi\right)_{(x)} \\
    & = \frac{1}{2}\left(\sum_{i<j} b_j \cdot b_i \cdot R^{\Sigma}\left(b_j, b_i\right) \phi + \sum_{i<j} b_i \cdot b_j \cdot R^{\Sigma}\left(b_i, b_j\right) \phi\right)_{(x)} \\
    & =- \frac{1}{2}\left(\sum_{i,j}^n b_j \cdot b_i \cdot R^{\Sigma}\left(b_i, b_j\right) \phi \right)_{(x)} 
  \end{aligned}$$
  wobei im letzten Schritt die Diagonalterme wegen Antisymmetrie der Krümmung verschwinden. Jetzt verwenden die Identität, welche die Krümmung des Spinorzusammenhangs mit der Ricci-Krümmung des Levi-Civita Zusammenhangs in Relation bringt:
  $$\begin{aligned} 
    &=-\frac{1}{4} \sum_{j=1}^n b_j \cdot \operatorname{Ric}\left(b_j\right) \cdot \phi \\ 
    & =-\frac{1}{4} \sum_{j=1}^n b_j \cdot \lambda_j b_j \cdot \phi=\frac{1}{4} \sum_{j=1}^n \lambda_j \cdot \phi \\ 
    & =\frac{1}{4} \operatorname{tr}(\operatorname{Ric}) \cdot \phi=\frac{1}{4} \mathrm{scal} \cdot \phi.
  \end{aligned}$$

\end{proof}


\begin{THEO}{}{Vanishing Theorem für harmonische Spinoren}
Sei $(M^n,g)$ eine kompakte Riemannsche Spin-Mannigfaltigkeit ohne Rand mit überall
positiver Skalarkrümmung $\operatorname{scal}_g>0$ und $D$ der
klassische Dirac Operator. Dann besitzt $D$ keine nichttrivialen
harmonischen Spinoren, d.\,h.
\[
\ker(D)=\{0\}.
\]
\end{THEO}

\begin{proof}
Sei $\psi \in \Gamma(\Sigma M)$ eine glatte Spinorsektion mit $D \psi=0$, dann gilt auch $D^2\psi=0$. Aus der Lichnerowicz-Formel folgt punktweise
$$
D^2 \psi=\nabla^{\Sigma*} \nabla^\Sigma \psi+\frac{1}{4} \operatorname{scal}_g \psi =0.
$$
Dann gilt
$$
0=\left( D^2 \psi, \psi\right)_{\Sigma M}=\left(\nabla^* \nabla \psi, \psi\right)_{\Sigma M}+\frac{1}{4} \int_M \operatorname{scal}_g\langle\psi,\psi \rangle_{\Sigma M} d \operatorname{vol}_g,
$$
und durch die formale Selbstadjungiertheit des Dirac-Operators gilt für den ersten Term:
$$
\left(\nabla^* \nabla \psi, \psi\right)_{\Sigma M}=(\nabla \psi, \nabla \psi)_{\Sigma M}=\|\nabla \psi\|_{\Sigma M}^2 .
$$
Damit erhalten wir
$$
0=\|\nabla \psi\|_{\Sigma M}^2+\frac{1}{4} \int_M \operatorname{scal}_g|\psi|^2_{\Sigma M} d \operatorname{vol}_g .
$$
Aus $\text {scal}_g>0$ überall auf $M$ folgt notwendigerweise:
$$
\|\nabla \psi\|_{\Sigma M}^2=0 \text { und } \int_M \operatorname{scal}_g|\psi|_{\Sigma M}^2 d \operatorname{vol}_g=0
$$
und damit
$$
\psi \equiv 0 .
$$
Also ist der Kern – und damit der Lösungsraum – des Dirac Operators trivial:
\[
\ker D=\{0\}.
\]
\end{proof}



\pagebreak

\section{Appendix}

\subsection{Algebren}

\begin{DEF}{}{Unitale $\mathbb{K}$-Algebra}
Eine \emph{unitale $\mathbb{K}$-Algebra} ist ein $\mathbb{K}$-Vektorraum $A$ zusammen mit
einer bilinearen Multiplikation
\[
A\times A\to A,\qquad (a,b)\mapsto ab,
\]
so dass $A$ mit dieser Multiplikation ein assoziativer Ring ist und ein Einselement
\[
1_A\in A
\]
besitzt, das heißt
\[
1_Aa=a1_A=a
\qquad\text{für alle }a\in A.
\]
\end{DEF}

\begin{DEF}{}{Algebrenhomomorphismus}
Seien $A$ und $B$ unitale $\mathbb{K}$-Algebren. Eine Abbildung
\[
\varphi:A\to B
\]
heißt \emph{$\mathbb{K}$-Algebrenhomomorphismus}, falls gilt:
\begin{enumerate}
    \item \(\varphi\) ist \(\mathbb{K}\)-linear,
    \item \(\varphi\) ist multiplikativ, d.h.
    \[
    \varphi(ab)=\varphi(a)\varphi(b)
    \qquad\text{für alle }a,b\in A,
    \]
    \item \(\varphi\) erhält das Einselement, d.h.
    \[
    \varphi(1_A)=1_B.
    \]
\end{enumerate}
Ist \(\varphi\) bijektiv, so heißt \(\varphi\) ein \emph{Algebrenisomorphismus}.
\end{DEF}

\subsection{Bündeltheorie}

\begin{DEF}{}{Faserbündel}
Seien $E$, $M$ und $F$ glatte Mannigfaltigkeiten. Ein \emph{Faserbündel} mit Totalraum $E$, Basisraum $M$ und typischer Faser $F$ ist eine surjektive glatte Abbildung
\[
\pi:E\to M,
\]
so dass es zu jedem $x\in M$ eine offene Umgebung $U\subset M$ und einen Diffeomorphismus
\[
\Phi_U:\pi^{-1}(U)\to U\times F
\]
gibt mit kommutierendem Diagramm
\[
\begin{tikzcd}
\pi^{-1}(U) \arrow[r,"\Phi_U"] \arrow[d,"\pi"'] & U\times F \arrow[dl,"\operatorname{pr}_1"] \\
U &
\end{tikzcd}
\]
Insbesondere ist für jedes $y\in U$ die \emph{Faser} $E_y:=\pi^{-1}(y)$ \emph{über} $y$ isomorph zu $\{y\}\times F\cong F$.
\end{DEF}

\begin{DEF}{}{Vektorbündel}
Ein \emph{Vektorbündel} ist ein Faserbündel $\pi:E\to M$ mit typischer Faser $V$ einem endlichdimensionalen Vektorraum $V$, sodass gilt:
\begin{enumerate}
  \item Jede Faser $E_x=\pi^{-1}(x)$ trägt die Struktur eines Vektorraums.
  \item Die lokalen Trivialisierungen
        \[
        \Phi_U:\pi^{-1}(U)\to U\times V
        \]
        sind faserweise linear.
\end{enumerate}
\end{DEF}

\begin{DEF}{}{Prinzipalbündel}
Sei $G$ eine topologische Gruppe bzw. Lie-Gruppe. Ein \emph{Prinzipalbündel}
mit Strukturgruppe $G$ ist ein Faserbündel $\pi:P\to M$ zusammen mit einer rechten Wirkung
\[
R: P\times G\to P,\qquad (p,g)\mapsto p\cdot g,
\]
so dass gilt:
\begin{enumerate}
    \item die Rechtswirkung \emph{erhält die Faser unter} $\pi$, d.h. 
    \[
    \begin{tikzcd}
    P \arrow[r,"R_g"] \arrow[dr,"\pi"'] & P \arrow[d,"\pi"] \\
    & M
    \end{tikzcd}
    \]
    kommutiert für alle $g\in G$,
    \item die Wirkung ist \emph{frei}, d.\,h.
    \[
    p\cdot g=p \;\Rightarrow\; g=e,
    \]
    \item die Wirkung ist auf jeder Faser $P_x:=\pi^{-1}(x)$ \emph{transitiv}, d.h. für alle $p,q\in P_x$ existiert ein $g\in G$, sodass 
    $$p\cdot g=q$$
    \item lokale Trivialisierungen $\Phi_u:\pi^{-1}(U)\to U\times G$ sind $G$-\emph{äquivariant}, d.h. das Diagramm
    \[
    \begin{tikzcd}
    \pi^{-1}(U)\times G \arrow[r,"\ \ \cdot\ "] \arrow[d,"\Phi_U\times \id_G"'] & \pi^{-1}(U) \arrow[d,"\Phi_U"] \\
    (U\times G)\times G \arrow[r,"\ \ \cdot\ "] & U\times G
    \end{tikzcd}
    \]
    kommutiert bzw. es gilt $(x,h)\cdot g=(x,hg)$.
\end{enumerate}
Insbesondere ist jede Faser $P_x$ ein \emph{freier transitiver $G$-Raum}.
\end{DEF}


\begin{DEF}{}{Assoziiertes Vektorbündel}
Sei \(\pi:P\to M\) ein Prinzipalbündel mit Strukturgruppe \(G\) und $\rho:G\to \GL(V)$ eine lineare Darstellung auf einem Vektorraum \(V\). Das \emph{assoziierte Vektorbündel}
ist der Quotient
\[
P\times_\rho V:=(P\times V)/\sim,
\]
wobei
\[
(p,v)\sim (p\cdot g,\rho(g^{-1})v).
\]
Die Projektion
\[
[p,v]\mapsto \pi(p)
\]
macht \(P\times_\rho V\) zu einem Vektorbündel über \(M\) mit typischer Faser \(V\) und Vektorraumstruktur auf jeder Faser durch
\[
[p,v]+[p,w]:=[p,v+w],
\qquad
\lambda [p,v]:=[p,\lambda v].
\]
\end{DEF}


\begin{DEF}{}{Zusammenhangsform auf $G$-Prinzipalbündeln}
  Sei $\pi:P\to M$ ein $G$-Prinzipalbündel. Eine Zusammenhangsform auf einem $G$-Prinzipalbündel ist eine $\mathfrak{g}$-wertige 1-Form
  $$
  \omega \in \Omega^1(P, \mathfrak{g}),
  $$
  für die gilt:
  \begin{enumerate}
    \item Reproduktion der Fundamentalfelder, d.h. 
    $$
    \omega\left(\xi^{\#}\right)=\xi \quad \forall \xi \in \mathfrak{g} .
    $$
    \item Äquivarianz unter Rechtswirkung
    $$
    \left(R_g\right)^* \omega=\operatorname{Ad}\left(g^{-1}\right) \omega \quad \forall g \in G .
    $$
    Punktweise bedeutet das:
    $$
    \omega_{p \cdot g}\left(\left(R_g\right)_* v\right)=\operatorname{Ad}\left(g^{-1}\right)\left(\omega_p(v)\right) \quad \forall v \in T_p P .
    $$
  \end{enumerate}
\end{DEF}



\subsection{Differentialoperatoren zwischen Vektorbündeln}



\begin{DEF}{}{Differentialoperator auf einem Vektorbündel der Ordnung $\le k$}
Sei $E \to M$ ein $\mathbb{K}$-Vektorbündel über einer glatten Mannigfaltigkeit $M$,
wobei $\mathbb{K}=\R$ oder $\C$.
Ein \emph{Differentialoperator der Ordnung $\le k$} auf $E$ ist eine lineare Abbildung
\[
P:\Gamma(E)\to \Gamma(E),
\]
so dass für jedes Koordinatensystem $x^1,\dots,x^n$ auf einer offenen Menge
$U\subset M$ und jede lokale Trivialisierung
\[
E|_U \cong U\times \mathbb{K}^p
\]
glatte Abbildungen
\[
A^\alpha:U\to \Mat(p\times p,\mathbb{K}),
\qquad |\alpha|\le k,
\]
existieren mit
\[
(Pv)|_U
=
\sum_{|\alpha|\le k}
A^\alpha(x)\,
\frac{\partial^{|\alpha|}v}{\partial (x^1)^{\alpha_1}\cdots \partial (x^n)^{\alpha_n}}
\]
für alle $v\in \Gamma(E)$.

Dabei ist
\[
\alpha=(\alpha_1,\dots,\alpha_n)\in \mathbb N_0^n,
\qquad
|\alpha|=\alpha_1+\dots+\alpha_n.
\]
\end{DEF}

\begin{PROP}{\parencite[p.~110]{BaerSpin}}{Zusammenhangslaplace}
  Sei $E$ ein Riemannsches oder Hermitesches Vektorbündel über einer Riemannschen Mannigfaltigkeit $M$ mit einem Metrik-Zusammenhang $\nabla$ auf E. Sei ferner $b_1,...,b_n$ ein lokaler ON-Rahmen in $TM$. Dann gilt für den Zusammenhangslaplace:
    $$
    \nabla^* \nabla=-\sum_{i=1}^n\left(\nabla_{b_i} \nabla_{b_i}-\nabla_{\nabla_{b_i}^{\mathrm{LC}}b_i}\right) .
    $$
\end{PROP}

\begin{DEF}{}{Prinzipalsymbol in Ordnung \(1\) und \(2\)}
Seien \(E,F\to M\) \(\mathbb{K}\)-Vektorbündel und \(x\in M\), \(\xi\in T_x^*M\).

\begin{enumerate}
    \item Ist \(P\in \mathscr D_{\le 1}(E,F)\), so heißt
    \[
    \sigma_1(P,\xi):E_x\to F_x,
    \qquad
    \sigma_1(P,\xi)e:=(P(f\widetilde e))(x),
    \]
    das \emph{Prinzipalsymbol} von \(P\), wobei \(f\in C^\infty(M)\) so gewählt ist, dass
    \[
    f(x)=0,\qquad df(x)=\xi,
    \]
    und \(\widetilde e\) eine glatte Fortsetzung von \(e\in E_x\) ist.

    Lokal, für
    \[
    P=\sum_{j=1}^n A^j(x)\partial_j + B(x),
    \]
    gilt
    \[
    \sigma_1(P,\xi)=\sum_{j=1}^n \xi_j A^j(x).
    \]

    \item Ist \(P\in \mathscr D_{\le 2}(E,F)\), so heißt
    \[
    \sigma_2(P,\xi):E_x\to F_x,
    \qquad
    \sigma_2(P,\xi)e:=\frac12\,(P(f^2\widetilde e))(x),
    \]
    das \emph{Prinzipalsymbol} von \(P\).

    Lokal, für
    \[
    P=\sum_{i,j=1}^n A^{ij}(x)\partial_i\partial_j + \sum_{j=1}^n B^j(x)\partial_j + C(x),
    \]
    gilt
    \[
    \sigma_2(P,\xi)=\sum_{i,j=1}^n \xi_i\xi_j A^{ij}(x).
    \]
\end{enumerate}
\end{DEF}

\begin{DEF}{}{Prinzipalsymbol als Bündelmorphismus}
Seien \(E,F\to M\) glatte \(\mathbb{K}\)-Vektorbündel und
\[
P:\Gamma(E)\to \Gamma(F)
\]
ein Differentialoperator der Ordnung \(\le k\). Das \emph{Prinzipalsymbol} von \(P\)
ist derjenige Schnitt
\[
\sigma_P\in \Gamma\bigl(\operatorname{Sym}^k(TM)\otimes \Hom(E,F)\bigr),
\]
dessen lokale Darstellung gerade durch den Anteil der höchsten Ableitungsordnung
von \(P\) gegeben ist.

Äquivalent dazu definiert \(\sigma_P\) für jedes \((x,\xi)\in T^*M\) eine lineare
Abbildung
\[
\sigma_P(x,\xi):E_x\to F_x,
\]
die homogen vom Grad \(k\) in \(\xi\) ist.
\end{DEF}



\pagebreak


\subsection{Motivation der Spin-Geometrie aus Mathematischer Physik}


Hier sollen drei Leitfragen adressiert werden: (a) Wieso man sich überhaupt für Spin-Geometrie interessieren sollte, (b) welche charakteristischen analytischen und algebraischen Eigenschaften den zu konstruierenden Dirac Operator auszeichnen und (c) welche zusätzlichen Anforderungen die globale Konstruktion eines solchen Differentialoperators an die zugrunde liegenden differentialgeometrischen Strukturen stellt. 



\subsubsection{Diracs Einführung des Dirac Operators in relativistischer Quantenmechanik}

In der ersten Hälfte des 20. Jahrhunderts war man nach der Entwicklung der Relativitätstheorie und der Quantenmechanik vor das Problem gestellt, beide Theorien zu vereinigen. Erste Schritte in der relativistischen Quantentheorie bestanden in der Suche nach Lorentzkovarianten Feldgleichungen auf dem Minkowski-Raum
$$\left(\mathbb{R}^{1,3}, \eta_{\mu \nu},(+,-,-,-)\right),$$
deren Lösungen eine Dynamik beschreiben, die mit der aus der speziellen Relativitätstheorie bekannten Energie-Impuls-Relation
$$E^2=|\mathbf{p}|^2+m^2$$
für freie Teilchen vereinbar ist.

Ein erster Ansatz lief über direkte Quantisierung (mit $\hbar=1, c=1$)
$$E\to i\partial_t \quad \mathbf{p}\to -i\nabla_x$$
woraus unmittelbar die EI-Relation erfüllt ist 
$$E^2\phi=(\mathbf{p}^2+m^2)\phi \quad \longleftrightarrow \quad -\partial_t^2\phi=(-\Delta_x+m^2)\phi$$
und mit $\square:=\partial_t^2-\Delta_x=\eta^{\mu\nu}\partial_\mu\partial_\nu$ zur Klein-Gordon Gleichung
$$(\square +m^2)\phi=0$$
führt. Diese wies jedoch Probleme konzeptioneller Natur auf (Erhaltungsstrom nicht positiv definit), und es wurde nach alternative Ansätze gesucht. 

Dirac verfolgte 1928 den Ansatz, eine relativistisch kovariante Feldgleichung zu entwickeln, welche a) aus einem linearen Differentialoperator 1. Ordnung entsteht, b) einen erhaltenen Strom mit positiv definiter Dichte besitzt (Interpretation als Wahrscheinlichkeitsdichte im Einteilchenbild) und c) die EI-Relation sicherstellt. 

Durch geeignete Faktorisierung des Klein-Gordon Operators wies Dirac die Existenz eines solchen Differentialoperators der Form
$$D=i\gamma^\mu\partial_\mu -m$$
mit den geforderten Eigenschaften nach, wobei der Operator $\slashed{\partial}:=i\gamma^\mu\partial_\mu$ als Dirac Operator bezeichnet wird. 

Dabei stellen sich in entscheidender Weise zwei wesentliche Unterschiede zur Klein-Gordon Gleichung heraus: Die Koeffizienten $\gamma^\mu$ des Dirac Operators müssen die Clifford Relation (Physikalische Konvention)
$$\gamma^\mu\gamma^\nu + \gamma^\nu \gamma^\mu= 2\eta^{\mu\nu}\mathbf{1}$$
erfüllen und sind damit Realisierungen eine Algebra-Darstellung $\rho: \Cl(\R^{1,3},\eta)\to \End(V)$ der Clifford Algebra $\Cl(\R^{1,3},\eta)$ in einem komplexen Vektorraum $V$. Dabei muss wegen $\mathrm{Cl}\left(\mathbb{R}^{1,3}\right) \otimes \mathbb{C} \cong \operatorname{Mat}(4, \mathbb{C})$ für nichttriviale Darstellungen die Dimension des komplexen Vektorraums $V$ mindestens $n\geq 4$ sein. Entsprechend wird der Dirac Operator im Gegensatz zur Klein-Gordon Gleichung nicht auf skalare Felder angewandt, sondern auf sogenannte Spinoren, als Funktionen der Form $\phi:\R^{1,3}\to V$. 

\subsubsection{Notwendigkeit der Clifford Relation für Dirac Operator}

Die Notwendigkeit der Clifford Relation werden beim Nachweis der EI-Relationen klar: Sei $\phi(x)=u(p)e^{-ip\cdot x}, p \cdot x:=p_\mu x^\mu$ mit $u(p)\neq 0$. Einsetzen in die Dirac Gleichung gibt
$$\left(i \gamma^\mu\left(-i p_\mu\right)-m\right) u(p) e^{-i p \cdot x}=0 \Longleftrightarrow\left(\gamma^\mu p_\mu-m\right) u(p)=0$$
Beide Seiten mit dem Faktor $\left(\gamma^\nu p_\nu+m\right)$ multiplizieren führt zu:
$$\left(\gamma^\nu p_\nu+m\right)\left(\gamma^\mu p_\mu-m\right) u(p)= (\gamma^\nu \gamma^\mu p_\nu p_\mu-m^2)u(p)=0$$
Unter der Annahme
$$\gamma^\nu \gamma^\mu p_\nu p_\mu \overset{(1)}{=} p^\mu p_\nu$$
erhalten wir also:
$$
\left(p^\mu p_\mu-m^2\right) u(p)=0 .
$$
und erzwingen zusammen mit der Identität $p^2=p^\mu p_\mu=E^2-|\mathbf{p}|^2$ und der Annahme $u(p)\neq 0$ die EI-Relation
$$E^2=|\mathbf{p}|^2+m^2$$
Um zu prüfen, unter welchen Bedingungen $(1)$ erfüllt wird, halten wir fest, dass $p_\nu p_\mu$ symmetrisch ist und damit nur der der symmetrische Teil von $\gamma^\nu \gamma^\mu$ zählt:
$$
\gamma^\nu \gamma^\mu p_\nu p_\mu=\frac{1}{2}\left(\gamma^\nu \gamma^\mu+\gamma^\mu \gamma^\nu\right) p_\nu p_\mu\overset{!}{=}\eta^{\nu \mu} p_\nu p_\mu \mathbf{1}=p^\mu p_\mu \mathbf{1}
$$
Wir finden also gerade physikalisch motiviert als notwendige Bedingung die Clifford Relation.\footnote{Alternativ könnte die Notwendigkeit der Clifford Relation gezeigt werden, indem man Diracs ursprüngliches Verfahren verfolgt, eine Faktorisierung des Klein-Gordon Operators von folgender Form zu finden. $(i\gamma^\mu\partial_\mu -m)(i\gamma^\mu\partial_\mu +m)=-\gamma^\mu\gamma^\nu\partial_\mu\partial_\nu -m^2=-(\square +m^2)$ via Symmetrie der partiellen Ableitungen und Clifford Relation.}


\subsubsection{Kovarianz der Dirac Gleichung bzgl $\Spin^+(1,3)$}

\begin{DEF}{}{Eigentliche orthochrone Lorentzgruppe}
Die eigentliche orthochrone Lorentzgruppe ist die Zusammenhangskomponente
der Identität in $\SO(1,3)$ und wird definiert durch
\[
\SO^{+}(1,3)
:=\{\Lambda\in\SO(1,3)\mid \det\Lambda=1,\ \Lambda^0{}_0>0\}.
\]
Sie erhält sowohl Raumorientierung als auch Zeitrichtung.
\end{DEF}

\begin{DEF}{}{Lorentz-Kovarianz der Dirac-Gleichung}
Die Dirac-Gleichung
\[
(i\gamma^\mu\partial_\mu - m)\psi = 0
\]
heißt \emph{Lorentz-kovariant}, wenn zu jeder
$\Lambda\in\SO^{+}(1,3)$ eine invertierbare lineare Abbildung
\[
S(\Lambda)\in\GL(V)
\]
existiert, so dass für
\[
\psi'(x):=S(\Lambda)\,\psi(\Lambda^{-1}x)
\]
die Äquivalenz
\[
(i\gamma^\mu\partial_\mu - m)\psi = 0
\quad\Longleftrightarrow\quad
(i\gamma^\mu\partial_\mu - m)\psi' = 0
\]
gilt.
\end{DEF}

\begin{REM}{}{Konsequenzen der Lorentz-Kovarianz}
Damit die Dirac-Gleichung lorentzkovariant ist, muss gelten:
\begin{enumerate}
\item \emph{Transformation der Clifford-Matrizen:}
\[
S(\Lambda)\,\gamma^\mu\,S(\Lambda)^{-1}
= \Lambda^\mu{}_\nu\,\gamma^\nu.
\]

\item \emph{Nicht-skalares Transformationsverhalten:}
Für skalare Felder wäre $S(\Lambda)=\id$, was die obige Relation unmöglich macht.

\item \emph{Nicht-tensorielle Natur:}
Tensorielle Darstellungen von $\SO^{+}(1,3)$ sind bei einer $2\pi$-Rotation trivial,
während Spinoren unter einer $2\pi$-Rotation das Vorzeichen wechseln.
\end{enumerate}
\end{REM}

\begin{LEM}{}{Fehlende echte Darstellung}
Es existiert keine echte lineare Darstellung
\[
S:\SO^{+}(1,3)\to\GL(V)
\]
mit
$$S(\Lambda) \gamma^\mu S(\Lambda)^{-1}=\Lambda_\nu^\mu \gamma^\nu$$
\end{LEM}

\begin{THEO}{}{Spingruppe als universelle Überlagerung}
Die universelle Überlagerung von $\SO^{+}(1,3)$ ist
\[
\lambda:\Spin^{+}(1,3)\to\SO^{+}(1,3),
\]
wobei $\Spin^{+}(1,3)\cong \SL(2,\mathbb{C})$ gilt.
\end{THEO}

\begin{THEO}{}{Spin-Darstellung und Clifford-Struktur}
Es existiert eine komplexe Darstellung
\[
\rho:\Spin^{+}(1,3)\to\GL(V)
\]
sowie eine Clifford-Multiplikation
\[
\gamma:\mathbb{R}^{1,3}\to\End(V)
\]
so dass für alle $g\in\Spin^{+}(1,3)$ und $v\in\mathbb{R}^{1,3}$
\[
\rho(g)\,\gamma(v)\,\rho(g)^{-1}
= \gamma(\lambda(g)v).
\]
gilt bzw. in Darstellung
$$\rho(g) \gamma^\mu \rho(g)^{-1}=\Lambda^\mu{ }_\nu \gamma^\nu$$
\end{THEO}

\begin{THEO}{}{Spin-Kovarianz der Dirac-Gleichung}
Die Dirac-Gleichung
$$
\left(i \gamma^\mu \partial_\mu-m\right) \psi=0
$$
ist damit in folgendem Sinne Spin-Kovariant: Es existiert eine Darstellung
$$
\rho:  \operatorname{Spin}^{+}(1,3) \rightarrow  \mathrm{GL}(V)
$$
und die Überlagerungsabbildung
$$
\lambda: \operatorname{Spin}^{+}(1,3) \rightarrow \mathrm{SO}^{+}(1,3),
$$
sodass für alle $g \in \operatorname{Spin}^{+}(1,3)$ gilt: 
\begin{enumerate}
  \item Transformation der Felder:
  $$
  \psi^{\prime}(x)=\rho(g) \psi\left(\lambda(g)^{-1} x\right)
  $$
  \item Kompatibilität mit Clifford-Multiplikation:
  $$
  \rho(g) \gamma(v) \rho(g)^{-1}=\gamma(\lambda(g) v) \quad \forall v \in \mathbb{R}^{1,3}
  $$
  \item Invarianz der Gleichung:
  $$
  \left(i \gamma^\mu \partial_\mu-m\right) \psi=0 \Longleftrightarrow\left(i \gamma^\mu \partial_\mu-m\right) \psi^{\prime}=0 .
  $$
\end{enumerate}
\end{THEO}

\begin{proof}

Sei $\Lambda=\lambda(g)$. Ableitung der Transformation  
$$
x \mapsto \Lambda x, \quad \psi(x) \mapsto \psi^{\prime}(x):=\rho(g) \psi\left(\Lambda^{-1} x\right) .
$$
gibt mit Kettenregel
$$
\partial_\mu \psi^{\prime}(x)=\rho(g)\left(\Lambda^{-1}\right)^\nu{ }_\mu\left(\partial_\nu \psi\right)\left(\Lambda^{-1} x\right) .
$$
Damit folgt aus der Transformationsregel für Clifford-Matrizen
$$
\gamma^\mu=\rho(g)^{-1} \Lambda^\mu{ }_\nu \gamma^\nu \rho(g) .
$$
Kontrahiere mit $\left(\Lambda^{-1}\right)^\sigma{ }_\mu$ :
$$
\left(\Lambda^{-1}\right)^\sigma{ }_\mu \gamma^\mu=\rho(g)^{-1}\left(\Lambda^{-1}\right)^\sigma{ }_\mu \Lambda^\mu{ }_\nu \gamma^\nu \rho(g)=\rho(g)^{-1} \gamma^\sigma \rho(g),
$$
weil
$$
\left(\Lambda^{-1}\right)^\sigma{ }_\mu \Lambda^\mu{ }_\nu=\delta^\sigma{ }_\nu .
$$
und damit
$$
\gamma^\mu\left(\Lambda^{-1}\right)_\mu^\nu=\rho(g)^{-1} \gamma^\nu \rho(g)
$$
Diese Formel verwenden wir:
$$
\begin{aligned}
\left(i \gamma^\mu \partial_\mu-m\right) \psi^{\prime}(x) & =i \gamma^\mu \rho(g)\left(\Lambda^{-1}\right)_\mu^\nu\left(\partial_\nu \psi\right)\left(\Lambda^{-1} x\right)-m \rho(g) \psi\left(\Lambda^{-1} x\right) \\
& =\rho(g)\left(i \gamma^\nu \partial_\nu-m\right) \psi\left(\Lambda^{-1} x\right) .
\end{aligned}
$$
Also gilt
$$
\left(i \gamma^\mu \partial_\mu-m\right) \psi=0 \quad \Longleftrightarrow \quad\left(i \gamma^\mu \partial_\mu-m\right) \psi^{\prime}=0 .
$$
\end{proof}

\begin{DEF}{}{Dirac Gleichung im masselosen Fall}
  Die Dirac Gleichung reduziert sich für den masselosen Fall $m=0$ also zu $\slashed{\partial}\phi=i\gamma^\mu \partial_\mu \phi=0$ für $\phi:\R^{1,3}\to V$. Damit gilt $\slashed{\partial}^2=-\square$.
\end{DEF}


\subsubsection{Übergang zur Spin-Geometrie}

Wir wollen nun zur Spin Geometrie übergehen. Der zuvor diskutierte physikalisch motivierte Dirac Operatoren im Minkowski Raum $(\R^{1,3},\eta)$ kann auf natürliche Weise im Rahmen der PseudoRiemannsche Geometrie verallgemeinert werden. Nachfolgend wollen wir jedoch stattdessen die Verallgemeinerung in der Riemannschen Geometrie verfolgen. Die entsprechenden Theorien werden zwar formal analog konstruiert, aber weisen wesentliche Unterschiede auf:

\begin{center}
\renewcommand{\arraystretch}{1.25}
\hspace*{-1cm}
\begin{tabular}{|l|c|c|}
\hline
\textbf{Aspekt} 
& \textbf{Riemannsche Spin-Geometrie} 
& \textbf{PseudoRiemannsche Spin-Geometrie} \\
\hline
Metrik 
& positiv definit 
& indefinit (Signatur $(p,q)$) \\
\hline
Typische Signatur 
& $(n,0)$ 
& $(1,3)$ (physikalisch relevant) \\
\hline
Clifford-Relation (Konvention)
& $v\cdot w + w\cdot v = -2g(v,w)$ 
& $v\cdot w + w\cdot v = 2g(v,w)$ \\
\hline
Spin-Gruppe 
& $\mathrm{Spin}(n)$ (kompakt) 
& $\mathrm{Spin}(p,q)$ (nicht kompakt) \\
\hline
Spinor-Darstellungen 
& unitär 
& nicht unitär \\
\hline
Dirac-Operator 
& elliptisch 
& hyperbolisch (oder gemischt) \\
\hline
Quadrat des Dirac-Operators 
& $D^2 = \nabla^*\nabla + \tfrac14\mathrm{scal}$ 
& Wellengleichungstyp \\
\hline
Analytischer Charakter 
& Spektraltheorie, Fredholm 
& Cauchy-Problem, Zeitentwicklung \\
\hline
Spektrum 
& diskret (kompakt) 
& kontinuierlich \\
\hline
Selbstadjungiertheit 
& kanonisch bzgl.\ $L^2$ 
& kein positives Skalarprodukt \\
\hline
Geometrische Rolle der Spinoren 
& topologische / analytische Invarianten 
& dynamische Felder \\
\hline
Physikalische Interpretation 
& keine (rein geometrisch) 
& Fermionfelder \\
\hline
Zentrale Anwendungen 
& Indexsatz, Positivität, Seiberg--Witten 
& Quantenfeldtheorie, Relativistik \\
\hline
\end{tabular}
\end{center}
\vspace{0.8cm}



Um einen Begriff von Dirac Operator zu entwickeln, welcher sinnvoll global auf Riemannsche Mannigfaltigkeiten definiert werden kann, wollen wir uns zunächst einen Eindruck davon verschaffen, welche Charakteristiken ein solcher Begriff erfüllen sollte, um der im Anschluss folgenden, nicht trivialen Konstruktion einen anleitenden Rahmen zu geben. Als Modellfall schauen wir uns dazu das euklidische Pendant des physikalischen Dirac Operator im Minkowskiraum an:



\begin{DEF}{}{Dirac Operator im euklidischen Raum}
  Sein $(\R^n,\delta_{\mu\nu})$ der $n$-dimensionale euklidische Raum mit Standardbasis $(e_1,…,e_n)$, trivialem Tangentialbündel $TM\cong \R^{n}\times \R^n$ und der aufgrund der Trivialität des Tangentialbündels und der konstanten Metrik global kanonisch definierbaren Clifford Algebra $\Cl(T_p\R^n,\delta_{\mu\nu}|_p)=\Cl(\R^n,\delta_{\mu\nu})$ mit Clifford Relation 
  $$v\cdot w + w \cdot v =-2 \langle v,w \rangle \mathbf{1},\ v,w\in\R^n.$$
  Sei $\Sigma_n$ ein irreduzibler komplexer Modul der Clifford-Algebra
  $\C\mathrm{l}\left(\mathbb{R}^n, \delta\right):=\Cl(\R^n)\otimes \C$.
  Das zugehörige (triviale) Spinor-Bündel über dem euklidischen Raum $\mathbb{R}^n$ ist definiert durch
  $$
  \Sigma \mathbb{R}^n:=\mathbb{R}^n \times \Sigma_n,
  $$
  wobei seine glatten Schnitte kanonisch mit glatten $\Sigma_n$-wertigen Funktionen identifiziert werden können,
  $$
  \Gamma\left(\Sigma \mathbb{R}^n\right) \cong C^{\infty}\left(\mathbb{R}^n, \Sigma_n\right) .
  $$
  Die Modulstruktur von $\Sigma_n$ wird durch eine Darstellung der Clifford-Algebra
  $$
  \rho: \C\mathrm{l}\left(\mathbb{R}^n, \delta\right) \longrightarrow \operatorname{End}\left(\Sigma_n\right)
  $$
  gegeben und schreiben für $z \in \C\mathrm{l}\left(\mathbb{R}^n\right)$ und $\psi \in \Sigma_n$
  $$
  z \cdot \psi:=\rho(z)(\psi) .
  $$
  Insbesondere wirkt jeder Vektor $v \in \mathbb{R}^n \subset \C\mathrm{l}\left(\mathbb{R}^n\right)$ durch Clifford-Multiplikation auf Spinoren.
  Unter Verwendung der euklidischen Metrik identifizieren wir das Kotangentialbündel mit dem Tangentialbündel,
  $$
  T_x^* \mathbb{R}^n \cong T_x \mathbb{R}^n \cong \mathbb{R}^n .
  $$
  Damit erhalten wir eine wohldefinierte Clifford-Multiplikation auf dem Spinor-Bündel als Bündelabbildung
  $$
  c: T^* \mathbb{R}^n \otimes \Sigma \mathbb{R}^n \longrightarrow \Sigma \mathbb{R}^n,
  $$
  die faserweise durch
  $$
  v^\flat \otimes \psi \mapsto v \cdot \psi
  $$
  gegeben ist. Weiter bezeichnen wir mit
  $$
  \partial: \Gamma\left(\Sigma \mathbb{R}^n\right) \longrightarrow \Gamma\left(T^* \mathbb{R}^n \otimes \Sigma \mathbb{R}^n\right)
  $$
  die komponentenweise partielle Ableitung von $\Sigma_n$-wertigen Funktionen. Explizit gilt
  $$
  (\partial \phi)(x)=\sum_{j=1}^n d x^j \otimes \partial_{e_j} \phi(x),
  $$
  wobei $\left(e_1, \ldots, e_n\right)$ die Standardbasis von $\mathbb{R}^n$ bezeichnet.
  
  Der (euklidische) Dirac-Operator ist nun als der lineare Differentialoperator erster Ordnung definiert durch
  $$
  D:=c \circ \partial: \Gamma\left(\Sigma \mathbb{R}^n\right) \longrightarrow \Gamma\left(\Sigma \mathbb{R}^n\right) .
  $$
  Für ein Spinorfeld $\phi \in \Gamma\left(\Sigma \mathbb{R}^n\right)$ ergibt sich damit mit $c(dx^j\otimes \partial_{e_j}\phi(x))=e_j\cdot \partial_{e_j}\phi(x)$ die explizite Darstellung
  $$
  D \phi(x)=\sum_{j=1}^n e_j \cdot \partial_{e_j} \phi(x) .
  $$
  Wählt man eine Darstellung der Clifford-Multiplikation durch Endomorphismen 
  $$\gamma^j:=c\left(e_j^b\right)=c\left(e_j^\flat\otimes \cdot \right) \in  \operatorname{End}\left(\Sigma_n\right),$$
  so lässt sich der Dirac-Operator äquivalent in Matrixform schreiben als
  $$
  D \phi(x)=\sum_{j=1}^n \gamma^j \partial_j \phi(x),
  $$
  wobei die Matrizen $\gamma^j$ die Clifford-Relation
  $$
  \gamma^i \gamma^j+\gamma^j \gamma^i=-2 \delta^{i j} \mathrm{Id}
  $$
  erfüllen.
\end{DEF}


\begin{REM}{}{Charakterisierende Eigenschaften des euklidischen Dirac Operators}
  Wir können folgende charakterisierende Eigenschaften des euklidischen Dirac Operators festhalten:
  \begin{enumerate}
    \item \emph{Erste Ordnung}: Der euklidische Dirac Operator ist offensichtlich ein linearer Differentialoperator erster Ordnung.
    
    \item \emph{Kompatibel mit Clifford-Multiplikation}: 
    Der euklidische Dirac-Operator ist mit der Clifford-Struktur verträglich. Diese Verträglichkeit äußert sich darin, dass der Dirac-Operator eine Leibniz-Regel erfüllt, bei der die Ableitung einer skalaren Funktion über die Clifford-Multiplikation in den Operator eingeht:
    $$
    D(f \phi)=(df) \cdot \phi+f D \phi \quad \text { für } f \in C^{\infty}\left(\mathbb{R}^n\right) .
    $$
    Man berechnet direkt mit $(dx^j)^\sharp=e_j$:
    $$\begin{aligned} D(f \phi) & =\sum_{j=1}^n e_j \cdot \partial_{e_j}(f \phi) \\ & =\sum_{j=1}^n e_j \cdot\left(\left(\partial_{e_j} f\right) \phi+f \partial_{e_j} \phi\right) \\ & =\sum_{j=1}^n\left(\partial_{e_j} f\right) e_j \cdot \phi+f \sum_{j=1}^n e_j \cdot \partial_{e_j} \phi \\ & = (d f)\cdot \phi+f D \phi.\end{aligned}$$

    \item \emph{Formal Selbstadjungiert}: Der euklidische Dirac Operator ist formal selbstadjungiert, d.h. für alle kompakt getragenen Spinorfelder $\phi, \psi$ gilt die Integralrelation
    $$
    \langle D \phi, \psi\rangle_{L^2}=\langle\phi, D^*\psi\rangle_{L^2}
    $$
    für $D=D^*$. 

    \item \emph{Elliptischer Operator mit Clifford-Multiplikation als Prinzipalsymbol}: Das Prinzipalsymbol von $D$ ist gegeben durch die Clifford-Multiplikation.
    Konkret gilt für $\xi \in T_x^* \mathbb{R}^n$, $f\in \Gamma(\Sigma \R^n)$ mit $f(x)=0,df(x)=\xi$ :
    $$
    \sigma_1(D, \xi)\phi= D (f\phi)(x)= df(x) \cdot \phi(x)+ f(x)(D\phi)(x))= c(\xi\otimes\phi)=\xi^\sharp\cdot \phi,
    $$
    Damit hat der quadrierte Dirac Operator $D^2$ das Prinzipalsymbol
    $$
    \sigma_2(D^2, \xi)=\sigma_1(D, \xi)\circ \sigma_1(D, \xi) = \xi^\sharp\cdot\xi^\sharp\overset{CR}{=} -|\xi^\sharp|^2 \mathrm{Id}=-|\xi|^2\mathrm{Id},
    $$
    was qua Definition zeigt, dass $D^2$ ein Operator vom Lapace-Typ ist. Aus der formalen Selbstadjungiertheit folgt zusätzlich, dass auch $D^*D=DD^*=D^2$ Operatoren vom Lapace-Typ sind.

    \item \emph{Lokal Natürlich}: Der euklidische Dirac-Operator ist lokal natürlich, d.h. er ist invariant unter euklidischen Isometrien, hängt ausschließlich von der euklidischen Metrik und der Clifford-Struktur ab, und benötigt keine zusätzlichen Wahlentscheidungen.
  \end{enumerate}
\end{REM}


\bigskip

\printbibliography[heading=bibintoc]
\end{document}